\documentclass[11pt,a4paper]{article}
\usepackage[margin=26mm,headheight=15pt]{geometry}
\usepackage{fontspec}
\setmainfont{texgyrepagella-regular.otf}[BoldFont=texgyrepagella-bold.otf,ItalicFont=texgyrepagella-italic.otf,BoldItalicFont=texgyrepagella-bolditalic.otf]
\usepackage{amsmath,amssymb}
\usepackage{unicode-math}
\setmathfont{texgyrepagella-math.otf}
\usepackage{graphicx,booktabs,longtable,array,tabularx}
\usepackage{xcolor,enumitem,fancyhdr}
\usepackage[unicode,colorlinks=true,linkcolor=teal!60!black,citecolor=teal!60!black,urlcolor=teal!60!black,bookmarksnumbered=true]{hyperref}
\usepackage{bookmark}
\usepackage{xurl}
\graphicspath{{../figures/}}
\setlength{\emergencystretch}{2em}
\setlength{\parskip}{3pt}
\clubpenalty=10000\widowpenalty=10000\displaywidowpenalty=10000
\raggedbottom
\renewcommand{\topfraction}{.9}
\renewcommand{\bottomfraction}{.85}
\renewcommand{\textfraction}{.08}
\renewcommand{\floatpagefraction}{.85}
\setcounter{topnumber}{2}\setcounter{bottomnumber}{2}\setcounter{totalnumber}{3}
\setlist{itemsep=2pt,topsep=4pt}
\renewcommand{\arraystretch}{1.16}
\setlength{\tabcolsep}{4pt}
\newcolumntype{Y}{>{\raggedright\arraybackslash}X}
\newcolumntype{L}[1]{>{\raggedright\arraybackslash}p{#1}}
\pagestyle{fancy}\fancyhf{}\fancyhead[L]{\small Three-Way reciprocal geometry}\fancyhead[R]{\small Reconstruction v0.1}\fancyfoot[C]{\thepage}
\newtheorem{theorem}{Theorem}[section]
\newtheorem{proposition}[theorem]{Proposition}
\newtheorem{lemma}[theorem]{Lemma}
\newtheorem{corollary}[theorem]{Corollary}
\newtheorem{definition}[theorem]{Definition}
\newtheorem{remark}[theorem]{Remark}
\newenvironment{proof}{\par\noindent\textit{Proof.}\ }{\hfill\(\square\)\par\medskip}
\numberwithin{equation}{section}
\newcommand{\PP}{\mathbb P}
\newcommand{\CC}{\mathbb C}
\newcommand{\RR}{\mathbb R}
\newcommand{\ZZ}{\mathbb Z}
\newcommand{\id}{\mathrm{id}}
\newcommand{\Gal}{\operatorname{Gal}}
\newcommand{\Aut}{\operatorname{Aut}}
\newcommand{\ord}{\operatorname{ord}}
\newcommand{\fig}[3]{\begin{figure}[htbp]\centering\includegraphics[width=\linewidth]{#1.pdf}\caption{#2}\label{#3}\end{figure}}
\hypersetup{pdftitle={Three-Way Reciprocal Geometry: Reconstruction of a Historical Rational Model},pdfauthor={Reconstruction draft prepared for Hilmir F. Halldorsson with Codex assistance},pdfsubject={Rational maps, branched covers, elliptic Galois closure; novelty-neutral draft}}
\begin{document}
\hypersetup{pageanchor=false}
\begin{titlepage}
\centering\vspace*{18mm}
{\Large Three-Way Reciprocal Geometry\par}\vspace{7mm}
{\LARGE\bfseries Reconstruction of a Historical Rational Model from Belyi Quotients to an Edwards Elliptic Galois Closure\par}
\vspace{15mm}
{\large Reconstruction draft prepared for Hilmir F. Halldórsson\par}
\vspace{3mm}{with Codex assistance\par}
\vspace{8mm}{Version 0.1 \quad 5 October 2026\par}
\vspace{12mm}
\begin{minipage}{.91\textwidth}
\textbf{Abstract.} We reconstruct the minus-sign reciprocal family
\[
F_{A,B}(y)=\frac{Ay^2-By+1}{y^2-By+A},\qquad y=x^2,
\]
from a historical visual construction and its later, partly inconsistent documentary record. Exact algebra separates locks, poles, cancellations, and real plotting restrictions. A coherent source normalization identifies a fixed degree-four Klein-four Belyi quotient with one marked-fiber parameter \(\lambda\). Retaining the original squaring cover introduces an independent parameter \(\nu\); for the selected covering type, the generic decorated coarse moduli space is \(M_{0,5}\). The degree-eight sphere map has a degree-sixteen genus-one Galois closure with group \(D_4\times C_2\). Its explicit equation is an Edwards curve, and \(\nu\) is a coordinate on \(X_0(4)\) for the closure with its selected cyclic subgroup of order four. The second marking specifies an effective fiber divisor, not a second elliptic modulus. We also prove the nested large-coefficient limits, explain the two distinct intermediate-chart involutions, and delimit an exact nearest-neighbor spectral-quotient comparison. The original real graph, the complex sphere cover, the elliptic closure, and any separately proposed surface of revolution are different objects. Historical physical claims are documented without being adopted.
\end{minipage}
\vfill
\begin{minipage}{.91\textwidth}\small
\textbf{Status.} This is a later mathematical reconstruction, not a revision of the historical papers. No mathematical novelty or priority is claimed. A dedicated prior-art audit of the assembled claims remains a separate stage. Author and publication metadata remain subject to author review; no submission or release is implied.
\end{minipage}
\end{titlepage}
\hypersetup{pageanchor=true}\pagenumbering{roman}
\tableofcontents\clearpage\pagenumbering{arabic}
\section{Purpose, scope, and the objects being reconstructed}\label{sec:intro}

Three-Way began as an exploratory comparison of reciprocal expressions, with particular attention to constants such as \(\pi,e,\phi\) and to the visible geometry of their graphs. That account of the original motivation is the author's retrospective account supplied with the present work order. It is not, by itself, a dated documentary demonstration of a particular formula. The surviving documents provide the more limited, checkable provenance developed in Section~\ref{sec:history}.

The construction accumulated physical language before its rational-map geometry had been classified. The present paper follows the opposite logical order: recover the formula, specify its domains, prove its properties, identify established mathematical structures, and only then delimit what a physical comparison actually establishes. Its subject is the original rational construction. A metaphor involving branes or universes is not an additional hypothesis in any theorem below.

Several distinct mathematical objects occur. A plotted real function is a subset of an affine real plane with specified axes. A rational map is a morphism between complex projective lines after common factors have been canceled. A decorated cover retains additional marked fibers or branch labels. A Galois closure is a larger function-field extension and can have a different source genus. Confusing these levels changes the question rather than answering it. In particular, a quartic-looking formula in \(x\) need not have genus one.

Our notation is
\begin{equation}\label{eq:family}
F(y)=F_{A,B}(y)=\frac{N(y)}{D(y)},\quad
N=Ay^2-By+1,\quad D=y^2-By+A,\quad f(x)=F(x^2).
\end{equation}
Parameters are complex unless a real statement is explicitly made. We use
\begin{equation}\label{eq:parameters}
c=A+1,\qquad \Delta=c^2-B^2,\qquad \chi=B^2-4A.
\end{equation}
The symbol \(f^{-1}_{\rm recip}\), when used historically, means \(1/f\), not a compositional inverse. The equation \(f=1/f\) is interpreted on the common finite, nonzero domain. Statements on \(\PP^1\) use the reduced rational map and local coordinates at infinity. An uncanceled numerical expression can retain holes that the projective map fills.

The reconstruction has three main classification levels. First, a nonconstant degree-two \(y\)-map sits inside a fixed degree-four \(V_4\) quotient; the marked quotient value is \(\lambda\). Second, retaining \(y=x^2\) gives a degree-eight map with additional branch value \(\nu\). Third, its generic normal closure is an elliptic curve with an order-sixteen group action. All of the named species---Belyi maps, Joukowski maps, marked spheres, Hurwitz covers, Edwards curves, and modular curves---are established mathematics. We prove their occurrence in the present family; we do not assign priority to that assembly.

\section{Formula provenance and documentary conflicts}\label{sec:history}

\subsection{Evidence hierarchy}
The historical PDFs and executable archive are primary evidence for what was printed or implemented~\cite{H1,H2,H3,H4}. The five external reconstruction reports R1--R5 are later analytical evidence, not historical sources. This paper supplies proofs independently of those reports. The source manifest records exact filenames, byte counts, SHA-256 hashes, and original locations. Verbatim snapshots are retained under \texttt{provenance}; no historical file has been altered.

Printed dates are reported as printed dates. Local modification times and a filename containing a date do not establish a publication date. The corpus associates H1 with Zenodo record 17240087 and H2 with record 17365334. Those associations are retained as corpus metadata; live record retrieval was not available during this draft's checks. We do not silently turn those associations into verified deposition timestamps. H3 and H4 are cited as the exact archived PDFs, without inventing a DOI.

\begin{table}[htbp]\centering\small
\begin{tabularx}{\linewidth}{L{12mm}L{26mm}Y}
\toprule ID & Printed date & Evidence and reconstruction consequence\\\midrule
H1 & 1 May 2025 & \emph{The Three-Way Quantum Gravitational Wave Model}, final draft. Appendix C, p.~13, prints minus-sign \((9,9)\) and \((\pi,e)\). Page 14 calls the integer example \((9,1)\) and prints \((\phi,1-\phi)\).\\
H2 & 16 October 2025 & \emph{Bounded Infinity and Resonant Curvature}. Pages 2--3 use the minus sign with a possible constant \(C\). Archived code defaults to \(C=1\). Labels and numerical CSVs do not consistently identify the integer pair.\\
H3 & 19 January 2026 & \emph{Experimental Geometry of a Three-Way Reciprocal Construct}. Page 3, Eq.~(1), explicitly uses \(+Bx^2\). Its numerical stability claims concern that convention.\\
H4 & 18 March 2026 & \emph{Reciprocal Rational Maps, Self-Dual Hinges, and Toroidal Phase Lifts}. Page 1 explicitly uses the plus sign. A phase lift is an additional construction.\\\bottomrule
\end{tabularx}
\caption{Principal documentary witnesses. These statements identify source content, not the truth of its interpretations.}\label{tab:provenance}
\end{table}

\subsection{The integer-pair discrepancy is substantive}
H1 prints
\begin{equation}\label{eq:historical99}
\frac{9x^4-9x^2+1}{x^4-9x^2+9},
\end{equation}
then calls an example \((9,1)\) on the following page. The archived H2 script \texttt{rrm\_tool.py}, in its preset routine, maps aliases including \texttt{9\_1} and \texttt{classic} to \((A,B,C)=(9,9,1)\), while its help text describes a classic \((9,1)\) example. Its \texttt{rrm\_asymptotes.csv} reports real poles near \(\pm1.0704662693\) and \(\pm2.8025170769\), exactly the \((9,9)\) regime. Conversely, \texttt{rtm\_key\_pairs.csv} explicitly has a row \(A=9,B=1,C=1\). These witnesses cannot all be reconciled by changing one printed digit.

We therefore retain both pairs as distinct documented examples. Their difference is mathematically large: \((9,9)\) has real zeros and poles; \((9,1)\) has none on the real \(x\)-axis. A figure or numerical output must be interpreted from its actual formula or recorded coefficients, not from the word ``classic.'' The executable preset and asymptote file corroborate the printed \((9,9)\) formula, but do not establish that every historical use of \((9,1)\) was a typo.

\subsection{The sign change is retained, not repaired}
The later plus-sign convention is
\[
F^+_{A,B}(y)=\frac{Ay^2+By+1}{y^2+By+A}=F_{A,-B}(y)=F_{A,B}(-y).
\]
As a family over complex parameters it is a reparametrization. At a fixed named pair with positive \(B\), however, it is a different real plot: replacing \(y\) by \(-y\) exchanges the positive and negative halves, and its \(x\)-lift is \(x\mapsto ix\), not a real reparametrization. The coherent minus-sign evidence in H1, H2, and the code archive, together with the explicitly selected reconstruction target, justifies Eq.~\eqref{eq:family}. It does not justify calling the later sign a transcription error.

H2 also writes \((Ay^2-By+C)/(y^2-By+A)\). This is a genuine extension. The numerator of the failure of reciprocity is
\[
N_C(y)y^2N_C(1/y)-D(y)y^2D(1/y)
=(C-1)\{A(y^4+1)-B(y^3+y)+(C+1)y^2\}.
\]
Thus the identity is preserved by \(C=1\), or by the exceptional case \(A=B=0,C=-1\), which gives \(-y^{-2}\); it does not hold for an arbitrary independently varied \(C\). The two-parameter reconstruction fixes \(C=1\).

\subsection{What the historical record does and does not establish}
The exact reciprocal identity and persistent visible intersections are recoverable mathematical content. H1's assertion about poles at the lock positions is not generally correct, and its suggestion that both \(+1\) and \(-1\) intersections occur in the displayed real examples requires correction. H2's finite sampled arc-length and curvature quantities depend on the interval, sampling, and treatment of poles; they are not by themselves invariant lengths of an unbounded singular graph. H3's perturbation and detectability results are historical numerical claims; the complete pipeline has not been independently reconstructed here. They are not used as proofs.

Modern exploratory files in the separate playground already derived several plus-sign identities and investigated constant choices. We record those as prior local analysis M1~\cite{M1}. Neither those notes nor the later reconstruction reports should be presented as evidence that the historical authors understood the subsequent Belyi or elliptic interpretation. The present paper is explicitly retrospective.

\section{Elementary algebra, locks, zeros, and poles}\label{sec:elementary}

\begin{proposition}[Reciprocity and level sets]\label{prop:levels}
After reduction, \(F(1/y)=1/F(y)\) and \(f(-x)=f(x)\). If \(A\ne1\) and \(\Delta\ne0\), then
\begin{align}
F=1&\iff y=1\ \text{or}\ y=-1,\label{eq:pluslevel}\\
F=-1&\iff c(y^2+1)-2By=0,\label{eq:minuslevel}
\end{align}
where the second equation is homogenized on \(\PP^1_y\). For \(c\ne0\), its roots are \((B\pm\sqrt{-\Delta})/c\); for \(c=0\), they are \(0,\infty\). The finite nonzero solutions of \(F=1/F\) are exactly these two level sets.
\end{proposition}
\begin{proof}
Reversing coefficients gives \(y^2N(1/y)=D(y)\) and \(y^2D(1/y)=N(y)\), so their quotient is reciprocal. Parity follows from \(y=x^2\). Subtraction and addition give
\begin{equation}\label{eq:sumdifference}
N-D=(A-1)(y^2-1),\qquad N+D=c(y^2+1)-2By.
\end{equation}
With no common root, these equations give the stated fibers, including their projective endpoints. Finally \(F=1/F\) is equivalent to \(F^2=1\) on its specified domain.
\end{proof}

Consequently the real locks \(x=\pm1\) are exact at height \(1\), provided \(A+1-B\ne0\). The other complex \(+1\) locks are \(x=\pm i\), requiring \(A+1+B\ne0\). They are equal-value intersections of two graphs, not necessarily fixed points of iteration \(x\mapsto f(x)\). Numerator/denominator reversal fixes them independently of the choice of named constants. The \(-1\) fiber is a different quadratic problem.

\begin{proposition}[Divisors and cancellations]\label{prop:divisors}
For \(A\ne0\), away from common factors, the zero and pole multisets of \(F\) are
\begin{equation}\label{eq:roots}
y_{z,\pm}=\frac{B\pm\sqrt\chi}{2A},\qquad
y_{p,\pm}=\frac{B\pm\sqrt\chi}{2}.
\end{equation}
The pole multiset is the reciprocal of the zero multiset, with the two displayed signs paired oppositely when square roots are coherently chosen. Common factors occur exactly on
\begin{equation}\label{eq:resultant}
\operatorname{Res}_y(N,D)=(A-1)^2\Delta=0.
\end{equation}
For \(A\ne\pm1\), the cancellation lines have reduced forms
\begin{align}
B=c:&\quad F(y)=\frac{Ay-1}{y-A},\label{eq:cancelplus}\\
B=-c:&\quad F(y)=\frac{Ay+1}{y+A}.\label{eq:cancelminus}
\end{align}
They have degree one. At \(A=1\), the reduced function is identically \(1\); at \((A,B)=(-1,0)\), it is identically \(-1\).
\end{proposition}
\begin{proof}
The quadratic formula gives \eqref{eq:roots}. A common zero for \(A\ne1\) must satisfy \(y^2=1\) by \eqref{eq:sumdifference}; substituting gives \(B=\pm c\). Expanding the resultant, or evaluating \(N-D\) at the roots of \(D\), gives \eqref{eq:resultant}. On \(B=c\), factor
\[
N=(y-1)(Ay-1),\qquad D=(y-1)(y-A).
\]
The negative line has the analogous factor \(y+1\). The two constant cases follow directly. No other degree drop is possible because the resultant is nonzero off these loci.
\end{proof}

The canceled point in \eqref{eq:cancelplus} has reduced value \(-1\), not \(+1\), when \(A\ne1\). Thus approaching a parameter cancellation and evaluating at the uncanceled lock need not commute. On \(B=-c\), the same statement holds at \(y=-1\). At \(A=1\), all apparent denominator roots are removable; the reduced map is constant, although the original expression leaves them undefined.

At \(A=0,B\ne0\), the projective zeros are \(1/B,\infty\) and the poles are \(B,0\), before the cancellations \(B=\pm1\). At \(A=B=0\), \(F=y^{-2}\). In \(x\), a nonzero simple \(y\)-root produces two simple roots; a root at \(y=0\) or infinity has its order doubled by squaring. Thus \(A=0\) is an embedding and marked-branch collision, even though the generic \(y\)-map still has degree two.

For a simple denominator zero \(x_0\) with nonzero numerator, the residue is
\[
\operatorname{Res}_{x=x_0}f=\frac{N(x_0^2)}{\left.\dfrac{d}{dx}D(x^2)\right|_{x=x_0}}.
\]
If numerator and denominator both vanish simply, their ratio tends to the ratio of derivatives and the singularity is removable. H1's pole argument conflated these alternatives. The ratio of derivatives is not the residue of a pole in that case.

\subsection{Critical points and parameter-independent statements}
Direct differentiation gives
\begin{equation}\label{eq:derivative}
F'(y)=\frac{(A-1)(-By^2+2cy-B)}{D(y)^2}.
\end{equation}
For a nonconstant degree-two map, the critical pair is
\[
y_{c,\pm}=\frac{c\pm\sqrt\Delta}{B}\quad(B\ne0),
\]
or \(0,\infty\) if \(B=0\). The points are reciprocal. If \(\chi=0\), one is a double zero and the other a double pole; a critical point at a pole is understood using the local target coordinate \(1/F\), not as an ordinary finite stationary graph point. For \(f\), the chain rule adds the ramification of \(x^2\) at \(0,\infty\).

Several visually persistent features now have exact explanations. Where defined, \(F(0)=1/A\) and \(F(\infty)=A\) are independent of \(B\). The product of the two nonzero finite pole roots is \(A\), and the zero product is \(1/A\), also independent of \(B\). The \(\pm1\) level locks are independent of both parameters off cancellation. The boundary \(\chi=0\) controls repeated zeros and poles; \(\Delta=0\) controls a different phenomenon, their cancellation. Sensitivities are
\begin{equation}\label{eq:sensitivity}
\partial_A F=\frac{(y^2-1)(y^2-By+1)}{D^2},\qquad
\partial_B F=\frac{(A-1)y(y^2-1)}{D^2}.
\end{equation}
Their common vanishing at the locks is algebraic. There is no distinguished integer \(9\) in these identities.

\subsection{Complete real-root classification}
For real parameters, exclude cancellation first. If \(A>0\), \(\chi<0\) gives nonreal conjugate \(y\)-roots; \(\chi>0\) gives real roots, both positive when \(B>0\), both negative when \(B<0\). At \(\chi=0\) the corresponding roots are double. If \(A<0\), each quadratic has one positive and one negative root. These claims follow from the discriminant, product, and sum. Taking both square roots of each \(y\)-root gives the full complex \(x\)-classification: positive roots give real pairs, negative roots imaginary pairs, and nonreal conjugate pairs give conjugate, parity-symmetric quartets.

\begin{table}[htbp]\centering\small
\begin{tabularx}{\linewidth}{L{49mm}Y}\toprule
Real coefficient regime & Real \(x\)-zeros and poles of the reduced nonconstant function\\\midrule
\(A>0,\chi<0\) & None.\\
\(A>0,\chi>0,B>0\) & Four simple zeros and four simple poles.\\
\(A>0,\chi>0,B<0\) & None; the corresponding roots are imaginary.\\
\(A>0,\chi=0,B>0\) & Two double zeros and two double poles.\\
\(A>0,\chi=0,B<0\) & None; double roots are imaginary.\\
\(A<0\) & Two simple real zeros and two simple real poles; also imaginary pairs.\\
\(A=0,B>0\) & Simple zeros \(\pm B^{-1/2}\), simple poles \(\pm\sqrt B\); double zero at infinity and double pole at zero.\\
\(A=0,B<0\) & Finite nonzero roots are imaginary; the orders at zero and infinity remain.\\
\(A=B=0\) & Pole order four at zero, zero order four at infinity.\\\bottomrule
\end{tabularx}\caption{Real plotting classification. The cancellation lines are handled separately by Proposition~\ref{prop:divisors}.}\label{tab:realroots}
\end{table}

For \(F=-1\), the two \(y\)-roots have product one. If \(\Delta>0\), they are nonreal conjugates of unit modulus. If \(\Delta<0\) and \(c\ne0\), they are real reciprocal roots and are positive exactly when \(B/c>0\), equivalently \(B/c>1\) in this regime. Only that case gives four real \(x\)-solutions. The case \(c=0\) gives the projective endpoints. Thus all four historical examples studied below, which have \(\Delta>0\), have no real \(-1\) branch, despite their exact real \(+1\) locks.

\section{The second involution and a coherent normal form}\label{sec:normal}

The reciprocal transformation exchanges the zero and pole fibers. A different transformation pairs points of equal value. It is already determined by the degree-two map, rather than being fitted to a graph.

\begin{theorem}[Klein-four normalization]\label{thm:normal}
Assume \(\Delta\ne0\). Define
\begin{equation}\label{eq:RT}
R(y)=1/y,\qquad T(y)=\frac{cy-B}{By-c}.
\end{equation}
These are distinct commuting Möbius involutions and generate \(V_4\). They satisfy \(F\circ R=1/F\) and \(F\circ T=F\). Choose \(d^2=\Delta\) and, coherently, put
\begin{equation}\label{eq:uk}
\rho=\frac{c+B}{d}=\frac{d}{c-B},\quad
u=\rho\frac{y-1}{y+1},\quad k=\frac{2(A-1)}d,
\quad r=u+u^{-1}.
\end{equation}
Then
\begin{equation}\label{eq:normal}
R:u\mapsto-u,\qquad T:u\mapsto u^{-1},\qquad
F=\frac{r+k}{r-k}.
\end{equation}
At \(A=1\) the last expression is interpreted as the constant rational map \(1\).
\end{theorem}
\begin{proof}
Cross-multiplying \(F(v)=F(y)\) and removing the diagonal factor gives
\[
N(v)D(y)-N(y)D(v)
=(A-1)(v-y)\{c(v+y)-B(vy+1)\}.
\]
For a nonconstant map, the other solution is \(v=T(y)\). The matrix
\(\left(\begin{smallmatrix}c&-B\\B&-c\end{smallmatrix}\right)\)
squares to \(\Delta I\) and has nonzero determinant \(-\Delta\), proving the involution assertion. Direct substitution gives commutation with \(R\). The coordinate \(u\) is Möbius because \(\rho\ne0\); substitution gives its two stated actions. Finally,
\[
r=\frac{2[c(y^2+1)-2By]}{d(y^2-1)}.
\]
Combining this with \eqref{eq:sumdifference} yields \eqref{eq:normal}.
\end{proof}

The two square roots must not be chosen independently. Replacing \(d\) by \(-d\) sends \((\rho,u,k,r)\) to \((-\rho,-u,-k,-r)\) and leaves \(F\) unchanged. Writing two unrelated principal square roots for \(\rho\) and \(k\) can introduce a false sign discrepancy in the normal form, particularly for complex parameters.

A generic orbit is \(\{u,-u,u^{-1},-u^{-1}\}\). The fixed pairs of \(R,T,RT\) are, respectively, \(\{0,\infty\}\), \(\{1,-1\}\), and \(\{i,-i\}\). Within the \(F\)-coordinate, \(T\) preserves value and \(R,RT\) reciprocate it. Thus the critical pair in Eq.~\eqref{eq:derivative} is exactly the fixed pair of the equal-value involution. This identifies its geometric role independently of a derivative calculation.

\fig{02_v4_orbit}{A generic Klein-four orbit in the normalized source plane. The orange pairing is reciprocity; the purple pairing is equal value. The unit circle is a coordinate guide, not the image of every real plot.}{fig:orbit}

\section{The fixed Belyi quotient and the reduced modulus}\label{sec:belyi}

\begin{theorem}[Universal quotient]\label{thm:belyi}
The quotient by the action in Theorem~\ref{thm:normal} is
\begin{equation}\label{eq:beta}
\beta(u)=\frac{(u^2+1)^2}{4u^2}=\frac{r^2}{4}.
\end{equation}
It has degree four, deck group \(V_4\), and ramification partitions \((2,2)\) over each of \(0,1,\infty\), with no other branch values. Moreover,
\begin{equation}\label{eq:marked}
\lambda=\frac{k^2}{4}=\frac{(A-1)^2}{\Delta},\qquad
\beta-\lambda=\frac{4ND}{\Delta(y^2-1)^2}.
\end{equation}
Thus the zero and pole sets form the marked quotient fiber \(\beta^{-1}(\lambda)\).
\end{theorem}
\begin{proof}
The function \(\beta\) is invariant under \(u\mapsto-u,1/u\). Its degree is four, so the four automorphisms account for the whole extension and
\(\CC(u)^{V_4}=\CC(\beta)\). It factors as the degree-two Joukowski map \(u\mapsto r=u+u^{-1}\) followed by \(r\mapsto r^2/4\). The fibers over zero and one are \(u=\pm i\) and \(u=\pm1\), each with multiplicity two; the poles \(0,\infty\) also have order two. Their total ramification is six, which equals \(2\cdot4-2\), so the list is complete. Squaring the expression for \(r\) in the preceding proof and using \((N+D)^2-(N-D)^2=4ND\) proves \eqref{eq:marked}.
\end{proof}

This is a standard Belyi map: a complex cover with branch values contained in \(\{0,1,\infty\}\), in the terminology reviewed by Sijsling and Voight~\cite{SV}. Its quotient orbifold is \(S^2(2,2,2)\), whose orbifold Euler characteristic is
\(2-3(1-1/2)=1/2\). Multiplying by the group order gives the source-sphere Euler characteristic two. The Klein-four orbifold is an established object~\cite{KO}. The original \(F\) has degree two in \(y\); it is not itself this degree-four map.

\fig{03_belyi_quotient}{The quotient factors through Joukowski and squaring maps. The bottom row identifies the three exceptional orbits. Branch labels are retained throughout the classification.}{fig:belyi}

\subsection{Exactly what equality of lambda means}
Retain the labeled quotient branch values \(0,1,\infty\), and mark the fiber containing zeros and poles; allow Möbius changes of source. In the generic reduced \(y\)-construction, equality of \(\lambda\) is necessary and sufficient for equivalence. Necessity follows because a labeled target isomorphism fixes the three normalized values and hence is the identity. For sufficiency, \(k'^2=k^2\), so \(k'=\sigma k\) for \(\sigma=\pm1\). The source change \(u'=\sigma u\) sends \(r'=\sigma r\) and preserves the full output value
\[
\frac{r'+k'}{r'-k'}=\frac{r+k}{r-k}.
\]
Composing with \(u\mapsto1/u\) gives the other equivalent choice. This also proves the assertion if the zero and pole sets are separately named: the isomorphism carries zero to zero and pole to pole between the two presentations.

The sign of \(k\) is needed in a fixed oriented \(u\)-chart. Changing it alone sends \(F\) to \(1/F\); changing \((u,k)\) simultaneously does not. Hence a signed coordinate choice is forgotten by \(\lambda\), but is not a further intrinsic modulus under the stated source equivalence. Labeling an individual zero, a sheet, or a fixed source coordinate is additional data not being classified here.

The lost embedding parameter can be recovered explicitly from \((\rho,k)\):
\begin{equation}\label{eq:recover}
A=\frac{\rho+\rho^{-1}+k}{\rho+\rho^{-1}-k},\qquad
B=\frac{2(\rho-\rho^{-1})}{\rho+\rho^{-1}-k}.
\end{equation}
The denominator must be nonzero for finite coefficients. These formulas follow by solving \(c/d=(\rho+\rho^{-1})/2\), \(B/d=(\rho-\rho^{-1})/2\), and \((A-1)/d=k/2\). They exhibit one marked-fiber coordinate and one source embedding coordinate, with a simultaneous sign ambiguity.

If both the branch labels and the zero/pole partition are forgotten, a still coarser equivalence identifies the six cross-ratio transforms
\[
\lambda,\ 1-\lambda,\ \lambda^{-1},\ (1-\lambda)^{-1},\
\lambda/(\lambda-1),\ (\lambda-1)/\lambda.
\]
For example, \(u\mapsto iu\) induces \(\beta\mapsto1-\beta\), and \(u\mapsto(u+1)/(u-1)\) induces \(\beta\mapsto\beta/(\beta-1)\). This coarser object is not the labeled object used in this paper.

\section{Real forms and the geometry of the plots}\label{sec:real}

For real coefficients with \(\Delta>0\), \(\rho,k\) may be chosen real. The real projective \(y\)-line maps to the real projective \(u\)-line. Since \(|u+u^{-1}|\ge2\) for nonzero real \(u\), its quotient image is \([1,\infty]\). For \(\Delta<0\), the same coherent choices are imaginary. Write \(u=iv\), \(v\) real. Then
\[
\beta=-\frac{(v-v^{-1})^2}{4}\le0.
\]
The real quotient image is the other projective arc, \((-\infty,0]\) with its endpoint at infinity. These are different real forms of the same complex quotient.

The identity
\begin{equation}\label{eq:chi-lambda}
\lambda-1=\frac{\chi}{\Delta}
\end{equation}
now explains the zero/pole regimes geometrically. If \(\Delta>0\), the marked level is positive. For \(0<\lambda<1\) it lies outside the real quotient image; for \(\lambda=1\) it meets a branch fiber and creates repeated roots; for \(\lambda>1\) it gives four real \(y\)-points. If \(\Delta<0\), then \(\lambda<0\) for a nonconstant family, and the marked fiber is real. Thus the statement ``\(\lambda<1\) means no real roots'' would be wrong without specifying the real form.

The real \(x\)-plot is stricter still: \(y=x^2\ge0\). In a real \(\rho\)-chart, this half-line maps to the segment joining \(-\rho\) and \(\rho\), not the whole real projective line. When \(\rho>1\), that segment contains \(u=\pm1\), hence the two nonzero real stationary points on the positive \(x\)-half-axis, provided they are finite graph values. If \(0<\rho<1\), it omits them. A real source Möbius equivalence at the \(y\)-level need not preserve the positive half-line, and need not lift to a real \(x\)-equivalence.

\fig{01_historical_graphs}{The four documented minus-sign examples, with \(f\) solid and \(1/f\) dashed. All share the exact \((\pm1,1)\) locks. Only \((9,9)\) has real zeros and poles. Curves are split at actual poles; vertical clipping is for display and does not join separate branches.}{fig:historical}
\fig{04_parameter_strata}{Real coefficient strata. The straight lines are cancellation, the parabola is repeated zeros/poles, and the dashed vertical line is the constant family. The dotted lines \(A=-1\) and \(B=0\) become distinguished in the full cover. The named examples occupy different marked-fiber and embedding regimes.}{fig:strata}

Complex conjugation, parity, source inversion, and output reciprocity have different roles. Real coefficients imply \(F(\overline y)=\overline{F(y)}\); parity acts in \(x\); \(x\mapsto1/x\) reciprocates the value. Reflecting a drawn graph about height one instead sends \(f\) to \(2-f\). That operation agrees with none of these globally for a nonconstant \(f\). Indeed, \(2-f=1/f\) implies \((f-1)^2=0\). It may be prescribed as a visual operation, but it is not an additional exact algebraic symmetry of this family.

\section{The diagonal large-coefficient family}\label{sec:scaling}

Let \(A=B=L\), with \(L\to+\infty\), and write \(f_L=F_{L,L}(x^2)\). The real transition at \(L=4\) is immediate from \(\chi=L(L-4)\). For \(0<L<4\), away from the constant case \(L=1\), zeros and poles are nonreal; at four they are double; for \(L>4\) they are real. At the threshold,
\[
F_{4,4}(y)=\frac{(2y-1)^2}{(y-2)^2}.
\]
The cancellation line intersects the positive diagonal nowhere; the excluded constant value is a separate event.

For \(L>4\), set
\[
q_L=\sqrt{\frac{1+\sqrt{1-4/L}}2}.
\]
The four positive distinguished roots, in increasing order, are
\begin{equation}\label{eq:Lroots}
z_{\rm in}=\frac1{\sqrt Lq_L}<z_{\rm near}=q_L<1<
p_{\rm near}=q_L^{-1}<p_{\rm out}=\sqrt Lq_L.
\end{equation}
Reciprocal pairing exchanges the first and last and the two middle points. Expansion gives \(q_L=1-(2L)^{-1}-5/(8L^2)+O(L^{-3})\). Thus a zero approaches the origin on scale \(L^{-1/2}\), a pole escapes on scale \(L^{1/2}\), and the two near-lock roots approach one on scale \(L^{-1}\). These are geometrically distinct uses of similar powers of \(L\).

\subsection{Bulk, inner, outer, and thin-lock limits}
The bulk identity is exact:
\begin{equation}\label{eq:bulk}
f_L(x)+x^2=\frac{x^6+1}{x^4+L(1-x^2)}.
\end{equation}
Consequently \(f_L\to-x^2\) locally uniformly on compact sets avoiding \(x=\pm1\). The reciprocal tends to \(-x^{-2}\) on compact sets also avoiding zero. At the locks, \(f_L(\pm1)=1\) for every \(L\); their limit cannot be read from the bulk formula. At zero the bulk absolute limit is correct, but it loses the leading value \(f_L(0)=L^{-1}\).

Three substitutions recover those missing scales:
\begin{align}
\zeta=Ly:&\quad
L F_{L,L}(\zeta/L)
=\frac{1-\zeta+\zeta^2/L}{1-\zeta/L+\zeta^2/L^3}
\longrightarrow1-\zeta,\label{eq:inner}\\
\eta=y/L:&\quad
\frac{F_{L,L}(L\eta)}L
=\frac{\eta^2-\eta/L+L^{-3}}{\eta^2-\eta+L^{-1}}
\longrightarrow\frac{\eta}{\eta-1},\label{eq:outer}\\
\tau=L(y-1):&\quad
F_{L,L}(1+\tau/L)
=\frac{1+\tau+\tau^2/L}{1-\tau+2\tau/L+\tau^2/L^2}
\longrightarrow\frac{1+\tau}{1-\tau}.\label{eq:thin}
\end{align}
Limits are locally uniform where the limiting denominators are nonzero. In \eqref{eq:outer}, the unreduced denominator vanishes at \(\eta=0\), so the stated derivation is uniform only away from \(0,1\); the simplified limiting expression's removable value at zero is not a claim of a uniform outer chart through the origin. The reciprocal profiles follow by inversion away from zeros: \(L^{-1}/F\to(1-\zeta)^{-1}\), \(L/F\to(\eta-1)/\eta\), and \(1/F\to(1-\tau)/(1+\tau)\), respectively.

The inner and bulk expansions overlap when \(1\ll|\zeta|\ll L\), where \((1-\zeta)/L\sim-y\). The outer and bulk expansions overlap when \(1\ll|y|\ll L\), equivalently \(L^{-1}\ll|\eta|\ll1\), where \(L\eta/(\eta-1)\sim-y\). Inversion exchanges the inner variable \(\zeta\) and outer variable \(\eta=1/\zeta\). In the thin chart it sends \(\tau\) exactly to \(-\tau/(1+\tau/L)\), tending to \(-\tau\), and reciprocates the limiting Möbius profile. These comparisons state explicit overlaps; no single truncated formula is asserted to be uniform across a pole.

\subsection{The intermediate shoulder profile, with matching}
Put \(h=L^{-1/2}\), \(s=(y-1)/h\), and \(P_h(s)=(f_L+1)/h\). Substituting into the exact rational expression gives
\begin{equation}\label{eq:Ph}
P_h(s)=\frac{s^2+2+2hs+h^2s^2}{-s+h(1+hs)^2}.
\end{equation}
\begin{theorem}[Intermediate limit and two-sided matching]\label{thm:matching}
On every compact subset of \(\CC\smallsetminus\{0\}\), for sufficiently small \(h>0\),
\begin{equation}\label{eq:W}
P_h(s)=W(s)-h(3+2/s^2)+O(h^2),\qquad
W(s)=-s-2/s.
\end{equation}
The convergence holds with any fixed number of derivatives. In the overlap
\(L^{-1}\ll|\delta|\ll1\), \(\delta=y-1\),
\begin{equation}\label{eq:twoterms}
f_L+1\sim-\left(\delta+\frac2{L\delta}\right)
\end{equation}
for real nonzero \(\delta\). The bulk-side overlap \(L^{-1/2}\ll|\delta|\ll1\) gives \(f_L+1\sim-\delta\); the lock-side overlap \(L^{-1}\ll|\delta|\ll L^{-1/2}\) gives \(f_L+1\sim-2/(L\delta)\).
\end{theorem}
\begin{proof}
On a compact set avoiding zero, the denominator in \eqref{eq:Ph} stays bounded away from zero for small \(h\). Taylor expansion in \(h\), including differentiation, gives \eqref{eq:W}; alternatively the exact difference is
\[
P_h-W=
\frac{h\{3s^2+2+h(3s^3+4s)+h^2(s^4+2s^2)\}}
{s[-s+h(1+hs)^2]}.
\]
To prove matching without an illicit interchange of limits, use the exact identity
\begin{equation}\label{eq:deltaexact}
f_L+1=-\frac{\delta+2/(L\delta)+(\delta+2)/L}
{1-(1+\delta)^2/(L\delta)}.
\end{equation}
For real \(0<|\delta|\le1/2\) and \(L|\delta|\ge5\), the relative difference from the two-term expression in \eqref{eq:twoterms} is at most \(10/(L|\delta|)\). Indeed the denominator correction is at most \(9/(4L|\delta|)\), hence its absolute denominator is at least \(11/20\); the extra numerator term relative to \(|\delta+2/(L\delta)|\) is at most \(5/(2L|\delta|)\). The triangle inequality gives the stated conservative bound. The two displayed overlap limits follow by comparing \(\delta^2\) with \(L^{-1}\).

For explicit comparison with the thin profile, set \(\tau=L\delta\). Its excess over \(-1\) is \(2/(1-\tau)\). Dividing the exact \(f_L+1\) by this expression gives
\[
\frac{[1+(\tau^2+2\tau)/(2L)+\tau^2/(2L^2)](1-\tau)}
{1-\tau+2\tau/L+\tau^2/L^2}=1+O(\tau^2/L)
\]
when \(1\ll|\tau|\ll\sqrt L\). Its large-\(|\tau|\) limit is \(-2/\tau\), exactly the lock-side intermediate term. This proves both matches in actual common domains.
\end{proof}

The real-relative statement avoids cancellation of the two leading terms. Over complex \(\delta\), those terms cancel at \(L\delta^2=-2\); the local analytic limit \eqref{eq:W} remains valid, but a uniform relative-error claim through those zeros would not be meaningful.

\subsection{Two involutions, two distinct finite-scale origins}
The source reciprocity \(y\mapsto1/y\) acts on the intermediate coordinate exactly as
\begin{equation}\label{eq:Rh}
R_h(s)=-\frac{s}{1+hs}\longrightarrow-s,
\qquad P_h(R_h(s))=-\frac{P_h(s)}{1-hP_h(s)}.
\end{equation}
Thus \(W(-s)=-W(s)\); reciprocity does not tend to \(s\mapsto2/s\). The equal-value involution from Theorem~\ref{thm:normal}, however, gives
\begin{equation}\label{eq:Jh}
\delta\mapsto\frac{\delta+2}{L\delta-1},\qquad
J_h(s)=\frac{hs+2}{s-h}\longrightarrow\frac2s,
\qquad P_h(J_h(s))=P_h(s).
\end{equation}
The two finite-
\(h\) Möbius transformations commute and square to the identity; they give an exact \(V_4\) action in this chart. Their limits generate \(s\mapsto-s,2/s\).

The same limiting equal-value symmetry exchanges the two matched terms in \eqref{eq:twoterms}: \(\delta\leftrightarrow2/(L\delta)\). It is therefore both visible as a balance symmetry of the matched expansion and inherited from the exact second involution. It is an emergent simple formula in the scaled coordinate, not a new independent symmetry of the original family, and specifically not the limit of reciprocal inversion.

Differentiating \eqref{eq:Ph} yields
\begin{equation}\label{eq:shoulders}
P_h'(s)=\frac{(1-h^2)(2+2hs-s^2)}{[-s+h(1+hs)^2]^2}.
\end{equation}
Its critical points are \(s=h\pm\sqrt{2+h^2}\), exactly the fixed points of \(J_h\). They tend to \(\pm\sqrt2\), the fixed points of \(s\mapsto2/s\), and
\(W'(s)=-1+2/s^2\) vanishes precisely there. The values are \(W(\pm\sqrt2)=\mp2\sqrt2\). These shoulders are the fold points of a degree-two equal-value quotient and the balance points of the two asymptotic contributions. This is a precise geometric interpretation of a visible feature, without an additional physical interpretation.

\fig{05_scaling_hierarchy}{The lock width and intermediate width are different powers of \(L\). On the right, the rescaled exact profiles approach \(W\), with stationary points approaching the purple marked points. The chart excludes \(s=0\); the exact lock lies in the thinner chart.}{fig:scales}

\subsection{How the quotient sees the charts}
For the diagonal family,
\[
\lambda_L=\frac{(L-1)^2}{2L+1},\qquad
\nu_L=\frac{(L+1)^2}{2L+1},\qquad
\nu_L-\lambda_L=\frac{4L}{2L+1}\longrightarrow2.
\]
Although both marked values diverge, their separation has a finite limit. Direct substitution in \eqref{eq:beta}, or in the marked-fiber identity, gives the chart descriptions
\begin{align*}
\text{bulk: }&\quad \beta/(L/2)\longrightarrow ((y-1)/(y+1))^2,\\
\text{thin lock: }&\quad \beta/(L/2)\longrightarrow\tau^{-2},\\
\text{intermediate: }&\quad u\longrightarrow s/\sqrt2,
\quad\beta\longrightarrow(s+2/s)^2/8,\\
\text{inner: }&\quad \beta-\lambda_L\longrightarrow2(1-\zeta),\\
\text{outer: }&\quad \beta-\lambda_L\longrightarrow2(1-1/\eta).
\end{align*}
Each statement excludes its displayed denominator zeros and any coordinate singularities; it is locally uniform on the remaining compact sets. Inversion pairs the inner and outer formulas. The equal-value involution sends fixed inner points toward \(\tau=-1\) and fixed outer points toward \(\tau=1\), collapsing further distinctions at leading thin-chart order. Resolving those images uniformly would require finer subcharts. We do not claim that this finite list is a compactification resolving every simultaneous limiting path. It is sufficient for the bulk, roots, lock, and shoulder regimes under discussion.

\section{Retaining the original squaring cover}\label{sec:xcover}

The reduced classification forgets where \(y=0,\infty\) lie on the normalized source. The original construction does not: those are the branch values of \(y=x^2\). Their images under the degree-four quotient coincide, introducing
\begin{equation}\label{eq:nu}
\nu=\frac{c^2}{\Delta}=\frac{(\rho+\rho^{-1})^2}{4},\qquad
\nu-1=\frac{B^2}{\Delta},\qquad \nu-\lambda=\frac{4A}{\Delta}.
\end{equation}
This is precisely the information lost when the extra cover is discarded.

\begin{theorem}[Degree-eight sphere cover]\label{thm:xcover}
If \(Bc\Delta\ne0\), the composite quotient is
\begin{equation}\label{eq:QAB}
Q_{A,B}(x)=\frac{[c(x^4+1)-2Bx^2]^2}{\Delta(x^4-1)^2}.
\end{equation}
It has degree eight and branch partitions \(2^4\) over \(0,1,\infty\), and \(2^2 1^4\) over \(\nu\). Its deck group is exactly \(\{x,-x,1/x,-1/x\}\simeq V_4\). In particular it is generically neither Galois nor a Belyi map after a target Möbius change.
\end{theorem}
\begin{proof}
Substitute \(y=x^2\) into \eqref{eq:beta}. Its numerator is nonzero at every root of \(x^4=1\), since \(c\pm B\ne0\), so its degree is eight. The identities
\begin{align}
Q-1&=\frac{[B(x^4+1)-2cx^2]^2}{\Delta(x^4-1)^2},\label{eq:Q1}\\
Q-\nu&=\frac{4x^2(c-Bx^2)(cx^2-B)}{\Delta(x^4-1)^2}\label{eq:Qnu}
\end{align}
give the stated multiplicities, with order two at both zero and infinity in the second line. Total ramification is \(4+4+4+2=14=2\cdot8-2\), proving completeness. The four stated source transformations preserve \(Q\).

For completeness of the deck group, a deck transformation preserves the uniquely identified ramified pair \(\{0,\infty\}\) over \(\nu\). It is therefore \(x\mapsto bx\) or \(b/x\). Preserving the pole set \(x^4=1\) forces \(b^4=1\), and preserving the nonzero mixed term of the numerator forces \(b^2=1\). Finally there are four distinct branch values, and the mixed partition at \(\nu\) is incompatible with a Galois cover, whose ramification indices over a fixed value must all agree.
\end{proof}

Writing \(t=B/c\) gives the simpler form
\begin{equation}\label{eq:Qt}
Q_t(x)=\frac{(x^4-2tx^2+1)^2}{(1-t^2)(x^4-1)^2},\qquad
\nu=\frac1{1-t^2}.
\end{equation}
For a fixed generic \(\nu\), the two signs of \(t\) give isomorphic complex labeled covers because \(Q_{-t}(ix)=Q_t(x)\). Conversely the moving labeled branch value is invariant. Thus \(\nu\) is a complete modulus for this selected family of covers with labeled branch values, before marking zeros and poles.

The second \(y\)-involution would lift rationally only if
\[
g(x)^2=\frac{cx^2-B}{Bx^2-c}.
\]
For \(Bc\Delta\ne0\), the right side has simple, distinct zeros and poles, so it is not a square in \(\CC(x)\). A rational square has even valuation at every point. The only nondegenerate exceptions are \(B=0\), where \(T(y)=-y\) lifts to \(ix\), and \(c=0\), where \(T(y)=-1/y\) lifts to \(i/x\). Their specialized maps are
\begin{equation}\label{eq:D4special}
Q_0(x)=\frac{(x^4+1)^2}{(x^4-1)^2},\qquad
Q_\infty(x)=-\frac{4x^4}{(x^4-1)^2}.
\end{equation}
Both are invariant under \(ix\) and \(1/x\), which generate a dihedral group of order eight, denoted \(D_4\) here. Since the map degree is eight, these are Galois covers. Their branch partitions are respectively \((2^4,4^2,2^4)\) and \((4^2,2^4,2^4)\) over \((0,1,\infty)\). They differ by \(Q_0=1-Q_\infty\), but that target relabeling is not allowed in the labeled problem.

\section{Two moduli, five marked points, and a selected Hurwitz type}\label{sec:moduli}

\subsection{The generic domain and the equivalence relation}
Let
\begin{equation}\label{eq:Omega}
\Omega=\{(\lambda,\nu):\lambda,\nu\notin\{0,1,\infty\},\ \lambda\ne\nu\}.
\end{equation}
An ordered five-pointed projective line is uniquely normalized by putting its first three labels at \(0,1,\infty\). Its remaining coordinates give exactly \(\Omega\), the familiar coordinate presentation of \(M_{0,5}\)~\cite[Eq.~(210)]{Mnev}. The statement about this target configuration is elementary; the statement that it classifies a cover requires additional proof and a specified covering type.

We retain the degree-two squaring tower, the three labeled quotient branches, the additional branch label \(\nu\), and the marked zero/pole fiber \(\lambda\). Individual sheets and individual zeros are not labeled. Source Möbius transformations are allowed; the target labels are preserved. Over the normalized fixed target, such a target transformation is the identity.

\begin{theorem}[Generic decorated classification]\label{thm:moduli}
For the Three-Way covering type, the generic decorated coarse moduli space is \(\Omega\simeq M_{0,5}\). Its coefficient domain is
\[
\mathcal P^\circ=\{(A,B)\in\CC^2:A(A-1)(A+1)B\Delta\chi\ne0\}.
\]
The map \((A,B)\mapsto(\lambda,\nu)\) is a four-sheeted finite étale cover of \(\Omega\), with fibers
\begin{equation}\label{eq:fourpairs}
(A,B),\quad(A,-B),\quad(A^{-1},B/A),\quad(A^{-1},-B/A).
\end{equation}
There is no additional algebraic constraint on \((\lambda,\nu)\) in this generic open set.
\end{theorem}
\begin{proof}
Introduce \(h_0=(A-1)/(A+1)\), distinct from the large-\(L\) variable \(h\). Then
\[
\lambda=\frac{h_0^2}{1-t^2},\qquad \nu=\frac1{1-t^2}.
\]
Given any point of \(\Omega\), choose the two independent square roots
\begin{equation}\label{eq:inversemoduli}
h_0^2=\lambda/\nu,\qquad t^2=1-1/\nu,
\qquad A=\frac{1+h_0}{1-h_0},\quad B=\frac{2t}{1-h_0}.
\end{equation}
Neither root vanishes, \(h_0\ne\pm1\), and \(t\ne\pm1\); all four choices give finite points in \(\mathcal P^\circ\). These are precisely \eqref{eq:fourpairs}. Algebraically, the two square roots are square roots of invertible regular functions on \(\Omega\), so they define a finite étale cover; the four choices are distinct. As a differential check,
\[
\det\frac{\partial(\lambda,\nu)}{\partial(A,B)}
=\frac{8B(A+1)(A-1)}{\Delta^3}\ne0.
\]

The coefficient transformations are realized in \(y\) by \(-y\) and \(1/y\), with complex \(x\)-lifts \(ix\) and \(1/x\). In particular
\(F_{A,-B}(-y)=F_{A,B}(y)\) and
\(F_{1/A,B/A}(1/y)=F_{A,B}(y)\).
Thus all four presentations define the same decorated object, even if zero and pole sets are separately named. Conversely equality of decorated objects preserves both target values. This proves completeness and surjectivity.
\end{proof}

The distinction between a marked whole fiber and separate zero/pole labels affects automorphisms, although not this coarse classification. For \(Q\), or \(Q\) with its entire regular fiber at \(\lambda\), the generic automorphism group is \(V_4\). Requiring the normalized \(F\), or its zero and pole sets separately, leaves only \(\{x,-x\}\simeq C_2\): inversion swaps the two sets. Five distinct labeled points on the target have trivial automorphism group. Therefore the covering stack is not the fine moduli space of five points; it retains stabilizers that the coarse space forgets.

\fig{07_five_marked_target}{The target configuration defining \(M_{0,5}\). The three fixed branch labels, the additional branch label, and the regular marked-fiber label have different covering roles. Their positions alone do not specify those roles or the covering type.}{fig:five}

\subsection{Why five points and the passport are not sufficient by themselves}
Over a regular \(\nu\), the universal quotient has the fiber
\(\{a,-a,a^{-1},-a^{-1}\}\), where here \(a\) is any one point of that fiber. A double cover of the source sphere branched at an unordered pair from this set again has a sphere as source. The six pairs fall into three \(V_4\)-orbits:
\[
\{a,-a\}\quad(R\text{-paired}),\qquad
\{a,a^{-1}\}\quad(T\text{-paired}),\qquad
\{a,-a^{-1}\}\quad(RT\text{-paired}).
\]
The Three-Way squaring construction selects the first type, since \(y=0,\infty\) map to \(u=-\rho,\rho\). The other two types give the same generic degree-eight ramification partitions but different labeled covers. Section~\ref{sec:galois} proves that the degree-four Galois intermediate is intrinsic and unique; any isomorphism of the whole cover must therefore preserve this pair-type distinction. Permuting target branch labels can relate the types, but our equivalence forbids that permutation.

This qualification locates the role of Hurwitz theory: \(M_{0,5}\) describes the selected component's coarse objects, not all covers sharing its target points and degree. No classification of unrelated monodromy types is needed here.

\subsection{Real presentations and fixed plotting coordinates}
For real coefficients in \(\mathcal P^\circ\), both square roots in \eqref{eq:inversemoduli} must be real. Hence the attainable real target region is exactly
\begin{equation}\label{eq:realmoduli}
(\nu>1,\lambda>0)\quad\text{or}\quad(\nu<0,\lambda<0),
\end{equation}
with the exclusions of \eqref{eq:Omega}. The converse follows from the inverse formulas. In particular the real coefficient family does not reach \(0<\nu<1\).

There are two real isomorphism classes among the four coefficient presentations over each attainable generic pair, distinguished by the sign of \(t\). Changing the sign of \(h_0\) is realized by the real transformation \(x\mapsto1/x\). Changing the sign of \(t\) requires \(ix\). To see that no alternative real Möbius map accomplishes it, any such map would preserve the ramified pair \(\{0,\infty\}\) and the pole set \(x^4=1\); it is therefore one of \(\pm x,\pm1/x\), all preserving \(t\). The actual graph retains even more: its fixed coordinate, chosen intervals, ordinate, and the restriction \(y\ge0\). A complex moduli point does not encode that graphical embedding.

\section{The genus-one normal closure and its monodromy}\label{sec:galois}

\begin{theorem}[Explicit Galois closure]\label{thm:galois}
For \(t\notin\{0,1,-1,\infty\}\), the smooth compactification in \(\PP^1_x\times\PP^1_z\) of
\begin{equation}\label{eq:Et}
E_t:\qquad tx^2z^2-x^2-z^2+t=0
\end{equation}
is the Galois closure of \(\CC(x)/\CC(Q_t)\). It has genus one and the map to the \(\beta\)-sphere has degree sixteen and group \(D_4\times C_2\), with \(|D_4|=8\).
\end{theorem}
\begin{proof}
Adjoin \(z\) by \(z^2=(x^2-t)/(tx^2-1)=T(x^2)\). The double cover of the \(x\)-sphere has four distinct branch points
\(\pm\sqrt t,\pm1/\sqrt t\). Riemann--Hurwitz gives
\(2g-2=2(-2)+4=0\), so its smooth compactification has genus one. The same branch calculation verifies smoothness of this normalized model; the displayed bidegree-\((2,2)\) compactification is smooth when those points are distinct.

Over \(\CC(y)\), the square classes of \(y\) and \((y-t)/(ty-1)\) are independent: their odd valuations occur respectively at \(0,\infty\) and \(t,1/t\), disjoint pairs. The degree over \(\CC(y)\) is four, hence sixteen over \(\CC(\beta)\). All conjugates of \(x\) are
\[
\pm x,\quad\pm1/x,\quad\pm z,\quad\pm1/z,
\]
so this field is the splitting field and therefore the normal closure.

The automorphisms
\begin{equation}\label{eq:generators}
\mathsf X(x,z)=(-x,z),\quad\mathsf Y(x,z)=(x,-z),\quad
\mathsf S(x,z)=(z,x),\quad\mathsf C(x,z)=(1/x,1/z)
\end{equation}
preserve the equation and \(\beta\). The first three generate the signed coordinate permutations, a group of order eight with \(\mathsf S\mathsf X\mathsf S=\mathsf Y\). The independent involution \(\mathsf C\) commutes with them. This gives sixteen distinct automorphisms, the whole extension degree, and the claimed direct product.
\end{proof}

\fig{06_cover_tower}{Degrees and genera of the reconstructed tower. The genus-one curve is an auxiliary normal closure of the degree-eight map. Passing to that closure changes the source; it does not change the genus of the original \(x\)-sphere.}{fig:tower}
\fig{08_elliptic_double_cover}{A cut-and-glue schematic for a real parameter \(0<t<1\), with \(a=\sqrt t\). Two sheets of the \(x\)-sphere, branched at four points, produce the compact genus-one surface. The drawing represents a complex double cover, not a surface obtained by rotating the real function graph.}{fig:closure}

\subsection{A reproducible branch-cycle description}
The degree-eight permutation action is on cosets of \(H=\langle\mathsf Y\rangle\). One branch-cycle tuple in loop order \(0,1,\nu,\infty\) is
\begin{align}
\sigma_0&=(1\,4)(2\,3)(5\,8)(6\,7),&
\sigma_1&=(1\,3)(2\,4)(5\,7)(6\,8),\nonumber\\
\sigma_\nu&=(1\,7)(2\,8),&
\sigma_\infty&=(1\,8)(2\,7)(3\,4)(5\,6).\label{eq:cycles}
\end{align}
With rightmost permutations acting first, their product is the identity. They are the coset actions of \(\mathsf S\mathsf C,\mathsf S,\mathsf X,\mathsf X\mathsf C\), respectively. They generate a transitive order-sixteen group; the cycle partitions agree with Theorem~\ref{thm:xcover}. Thus this tuple realizes the stated cover rather than merely listing arbitrary permutations with matching cycle types.

For an intrinsic intermediate-cover argument, the normal closure of \(H\) in the Galois group is
\(N=\langle\mathsf X,\mathsf Y\rangle\), of order four. Any degree-four intermediate extension that is Galois over \(\CC(\beta)\) corresponds to a normal order-four subgroup containing \(H\), which must contain its normal closure and hence equal \(N\). This proves uniqueness of the degree-four Galois intermediate and completes the pair-type distinction used in Section~\ref{sec:moduli}.

\section{Collision boundaries and specialized maps}\label{sec:boundary}

The identities \eqref{eq:chi-lambda} and \eqref{eq:nu} give a dictionary between coefficient strata and collisions of target labels.
\begin{table}[htbp]\centering\small
\begin{tabularx}{\linewidth}{L{28mm}L{27mm}Y}\toprule
Coefficient locus & Target event & Meaning away from intersections\\\midrule
\(A=1\) & \(\lambda=0\) & The nonconstant \(F\)-decoration disappears; \(Q\) can remain degree eight.\\
\(B^2=4A\) & \(\lambda=1\) & Double zeros and poles; the marked fiber meets a branch value.\\
\(A=0\) & \(\lambda=\nu\) & Zeros/poles meet the additional branching from squaring.\\
\(B=0\) & \(\nu=1\) & Dihedral degree-eight Galois specialization.\\
\(A=-1, B\ne0\) & \(\nu=0\) & The other labeled dihedral specialization.\\
\(\Delta=0\) & \(\lambda,\nu\to\infty\) & Generically a triple target collision and cancellation of \(F\).\\\bottomrule
\end{tabularx}\caption{Different exceptional loci have different geometric meanings. A marked-point collision need not lower the degree of the unmarked cover.}\label{tab:boundaries}
\end{table}

The standard stable-target compactification is
\begin{equation}\label{eq:barM}
\overline M_{0,5}=\operatorname{Bl}_{(0,0),(1,1),(\infty,\infty)}
(\PP^1_\lambda\times\PP^1_\nu).
\end{equation}
Its ten boundary divisors correspond to partitions of the five labels into a pair and a triple~\cite[\S2]{Braungardt}. Seven are the proper transforms of the six coordinate lines and the diagonal; the other three are exceptional divisors. A collision is represented by an additional rational target component carrying the colliding labels. The labels are retained, not deleted.

For the cover, the standard limit is an admissible cover, with nodes mapping to nodes and equal local ramification on the two branches. The local expressions have the form \(\xi\mapsto\xi^e\), \(\eta\mapsto\eta^e\), with the source and target smoothing parameters related accordingly~\cite[Definition 8; \S5.1]{CMR}. Products of colliding branch cycles in \eqref{eq:cycles} have types \(4^2\) when \(\nu\to0\) or one, and \(2^4\) when \(\nu\to\infty\). The regular marking \(\lambda\) has identity monodromy, so its collisions do not add a permutation factor.

At a generic point of \(\Delta=0\),
\[
\lambda/\nu=((A-1)/(A+1))^2
\]
retains the relative position of the colliding labels \(\infty,\lambda,\nu\). The stable target lies on the exceptional divisor over \((\infty,\infty)\), with partition \(\{0,1\}\mid\{\infty,\lambda,\nu\}\). Reducing the canceled rational expression \(F\) to degree one loses the other components and is not the same operation as taking this degree-preserving limit. At intersection points such as \((1,\pm2)\) or \((-1,0)\), approach directions and rates can matter; an indeterminate coefficient point alone does not choose a stable marked limit.

One precise coarse statement uses the normalization of the closure of the selected Hurwitz component in an admissible-cover compactification, without individual sheet labels. Its branch map is proper with finite fibers: over a fixed stable marked target, there are finitely many monodromy choices and compatible node gluings. It is therefore finite, and it is birational because the interior classification is one-to-one. A finite birational morphism to the normal surface \(\overline M_{0,5}\) is an isomorphism. This concerns the normalized coarse component only. It does not identify boundary stacks or a fixed-coordinate rational formula with the stable compactification.

Finally, the rational dihedral maps \eqref{eq:D4special} result from merging branch labels. Their rational Galois closures have degree eight. A degeneration of the generic degree-sixteen genus-one Galois cover, with all labels retained, is a different object with additional components. The two descriptions are compatible, but their source genera and degrees should not be interchanged.

\section{The closure as an Edwards elliptic curve}\label{sec:edwards}

Choose \(a_e^2=t\) and take the elliptic origin \(O=(0,a_e)\) on \(E_t\). The subscript distinguishes this square root from the source-coordinate scale \(\rho\) and from the generic orbit point used earlier. Choosing an origin does not add a modulus of the complex genus-one curve. There is also the section \((1,i)\) over \(\CC(t)\); thus the family has a point without adjoining \(\sqrt t\), although \(O\) makes the group formulas simpler on each fiber.

Set \(U=x/a_e\), \(V=z/a_e\). The equation becomes
\begin{equation}\label{eq:Edwards}
U^2+V^2=1+t^2U^2V^2,\qquad O=(0,1).
\end{equation}
This is exactly an untwisted Edwards form, a standard elliptic model~\cite{Edwards,Morain,EFD}. The excluded values \(t=0,\pm1,\infty\) are essential. We use its smooth compactification; the affine real curve and a singular plane-quartic closure are not substitutes for that compact curve.

\subsection{Montgomery and Weierstrass equations}
Define on a dense open set
\[
M=\frac{1+V}{1-V},\qquad N_M=\frac{M}{U}.
\]
Substituting \(U^2=(1-V^2)/(1-t^2V^2)\) gives
\begin{equation}\label{eq:Montgomery}
\gamma N_M^2=M^3+\alpha M^2+M,
\qquad \alpha=\frac{2(1+t^2)}{1-t^2}=4\nu-2,
\quad\gamma=\frac4{1-t^2}=4\nu.
\end{equation}
The inverse is \(U=M/N_M\), \(V=(M-1)/(M+1)\). Rational inverse maps between smooth projective curves extend over the finitely many omitted coordinate points, so these are isomorphisms of compact curves. The origin goes to infinity. Over \(\CC\), put \(W=\sqrt\gamma N_M\), obtaining
\begin{equation}\label{eq:Weierstrass}
W^2=M^3+(4\nu-2)M^2+M.
\end{equation}
The square root can affect a twist over a smaller field; it does not change the complex isomorphism class being classified.

For the short Weierstrass form put \(M=\xi-\alpha/3\). Then
\begin{align}
W^2&=\xi^3+p\xi+q,\nonumber\\
p&=-\frac{16\nu^2-16\nu+1}{3},\qquad
q=\frac{2(2\nu-1)(32\nu^2-32\nu-1)}{27},\label{eq:short}\\
\Delta_W&=-16(4p^3+27q^2)=256\nu(\nu-1).\label{eq:Wdisc}
\end{align}
The subscript on \(\Delta_W\) distinguishes it from the coefficient discriminant \(\Delta\). Its displayed form is an affine model on the finite \(\nu\)-line; analyzing infinity requires changing coordinates and, if desired, choosing a minimal model.

\subsection{Jacobi quartic and a fully explicit Legendre map}
Set \(w=(tx^2-1)z\). Squaring and using \eqref{eq:Et} gives
\begin{equation}\label{eq:quartic}
w^2=(tx^2-1)(x^2-t)=tx^4-(t^2+1)x^2+t.
\end{equation}
Thus \(v=w/a_e\) satisfies the Jacobi-quartic equation
\[
v^2=x^4-(t+t^{-1})x^2+1.
\]
This is another established model, with Jacobi parameter \(-\tfrac12(t+t^{-1})\) in the convention \(v^2=x^4+2a_Jx^2+1\)~\cite{EFD}. The inverse recovers \(z=w/(tx^2-1)\).

Define
\begin{equation}\label{eq:Legendrechange}
q_0=\frac{1-t}{1+t},\quad
\ell=q_0\frac{x-a_e}{x+a_e},\quad
\kappa=\frac{2ia_e(t-1)}{(t+1)^2},\quad
\eta_E=\kappa\frac{w}{(x+a_e)^2}.
\end{equation}
Then
\begin{equation}\label{eq:Legendre}
\eta_E^2=\ell(\ell-1)(\ell-m),\qquad
m=\left(\frac{t-1}{t+1}\right)^2.
\end{equation}
To verify the constant as well as the branch cross-ratio, the four branch points map as
\[
a_e\mapsto0,\quad -a_e\mapsto\infty,\quad
-a_e^{-1}\mapsto1,\quad a_e^{-1}\mapsto m,
\]
and direct multiplication yields
\[
\ell(\ell-1)(\ell-m)
=-\frac{4t(t-1)^2}{(t+1)^4}\frac{w^2}{(x+a_e)^4}.
\]
This is \(\kappa^2w^2/(x+a_e)^4\). The inverse starts with
\(x=a_e(q_0+\ell)/(q_0-\ell)\) and then recovers \(w,z\). This Legendre map uses \((-a_e,0)\) as the point sent to infinity, so it changes the elliptic origin relative to \(O\). It leaves the bare curve and its \(j\)-invariant unchanged. Reordering the four branch points gives the usual six cross-ratio transforms of \(m\).

\fig{09_edwards_legendre}{Affine real representatives for \(t=1/2\) and \(m=1/9\). The Edwards origin and order-four point are marked. These panels display real equations, not a homeomorphism between their plotted real loci: the displayed Legendre change uses an imaginary factor and is asserted over \(\CC\). Both compact complex curves have genus one.}{fig:models}

\subsection{Two independent computations of j}
For \eqref{eq:Weierstrass}, the standard Weierstrass invariants computed from its coefficients are
\[
c_4=16(\alpha^2-3),\quad\Delta_W=16(\alpha^2-4),\quad
j=256\frac{(\alpha^2-3)^3}{\alpha^2-4}.
\]
Substitution gives the two equivalent expressions
\begin{equation}\label{eq:jt}
j(t)=16\frac{(t^4+14t^2+1)^3}{t^2(t^2-1)^4},
\end{equation}
\begin{equation}\label{eq:jnu}
j(\nu)=16\frac{(16\nu^2-16\nu+1)^3}{\nu(\nu-1)}.
\end{equation}
Independently, the Legendre expression
\(256(1-m+m^2)^3/[m^2(1-m)^2]\), with \(m\) from \eqref{eq:Legendre}, reduces to \eqref{eq:jt}. These are exact algebraic identities, not numerical recognition of an elliptic invariant. The standard Edwards/Montgomery invariant framework is discussed in Morain~\cite[\S\S2.4--2.5]{Morain}; all substitutions needed here have been displayed.

\section{Translations, inversion, and the two Klein-four actions}\label{sec:group}

The abstract Galois group acquires a concrete geometric meaning on \(E_t\). On the Edwards equation, addition is
\[
(U_1,V_1)+(U_2,V_2)=
\left(\frac{U_1V_2+V_1U_2}{1+t^2U_1U_2V_1V_2},
\frac{V_1V_2-U_1U_2}{1-t^2U_1U_2V_1V_2}\right),
\]
where defined, extending on the smooth compact curve~\cite{EFD}. Negation is \((-U,V)\). Let \(P=(a_e,0)\) in the original coordinates. The formula gives
\[
2P=(0,-a_e),\qquad4P=O,\qquad
\tau_P(x,z)=(z,-x),\qquad\tau_{2P}(x,z)=(-x,-z).
\]
Thus \(P\) has exact order four and
\begin{equation}\label{eq:translationactions}
\mathsf X=[-1],\quad
\mathsf Y=\tau_{2P}[-1],\quad
\mathsf S=\tau_P[-1].
\end{equation}

The remaining involution \(\mathsf C\) has no fixed point on a smooth fiber. A fixed point would require \(x,z\in\{1,-1\}\), forcing \(t=1\). Every elliptic involution with linear part \(-1\) has a fixed point, since \([2]\) is surjective on a complex elliptic curve. A fixed-point-free involution must therefore be translation by a nonzero two-torsion point. Here it is
\[
Q_2=\mathsf C(O)=(\infty,a_e^{-1}),\qquad Q_2\ne2P.
\]
Consequently
\begin{equation}\label{eq:Gelliptic}
H_{\rm tr}=\langle P,Q_2\rangle\simeq C_4\times C_2,
\qquad G=H_{\rm tr}\rtimes\langle[-1]\rangle\simeq D_4\times C_2.
\end{equation}
Inversion negates the order-four generator and fixes the order-two generator. This identifies the order-four group element as a translation, rather than merely assigning it a dihedral name.

The full two-torsion is \(E_t[2]=\{O,2P,Q_2,2P+Q_2\}\), at the four points
\[
(0,\pm a_e),\qquad(\infty,\pm a_e^{-1}).
\]
The four branch points of the \(x\)-projection form the coset \(P+E_t[2]\), with coordinates \((\pm a_e,0)\) and \((\pm a_e^{-1},\infty)\). This connects the branch calculation directly to elliptic torsion.

The original source sphere is the Kummer quotient
\begin{equation}\label{eq:Kummerx}
\PP^1_x=E_t/\langle\mathsf Y\rangle.
\end{equation}
Since \(\mathsf Y(R)=-R+2P\), it is inversion after shifting the origin to \(P\). Two-torsion translations commute with it and descend faithfully to \(x\mapsto\pm x,\pm1/x\). This is the \(V_4\) deck action on the degree-eight sphere map.

For the other sphere,
\[
\PP^1_y=E_t/\langle\mathsf X,\mathsf Y\rangle.
\]
Let \(E'=E_t/\langle2P\rangle\). Then the \(y\)-sphere is the Kummer quotient of \(E'\), and
\(H_{\rm tr}/\langle2P\rangle=E'[2]\). Its two-torsion translations induce the original normalized \(u\)-action by \(-u,1/u\). The two Klein-four actions are thus related through an isogeny, but they act on different spheres and arise from different subgroups before quotienting.

\section{The modular meaning of nu and the natural two-isogeny}\label{sec:modular}

\begin{theorem}[The selected cyclic subgroup]\label{thm:X04}
The coordinate \(\nu\) is a Hauptmodul on \(X_0(4)\) for the pair \((E_t,\langle P\rangle)\). On the smooth locus,
\[
Y_0(4)\simeq\PP^1_\nu\smallsetminus\{0,1,\infty\}.
\]
Forgetting the cyclic subgroup gives the degree-six map \eqref{eq:jnu} to the \(j\)-line. In particular \(\nu\) retains more than the bare complex elliptic isomorphism class.
\end{theorem}
\begin{proof}
In the Weierstrass model, the order-four generator and its double are
\[
P=(1,\sqrt{\alpha+2}),\qquad2P=(0,0).
\]
Conversely take a complex elliptic curve with a cyclic order-four subgroup \(\langle P\rangle\). Put its nonzero double at \((0,0)\) in an equation \(Y^2=X(X^2+bX+c_0)\). The duplication formula gives
\[
X(2P)=\frac{(X(P)^2-c_0)^2}{4X(P)(X(P)^2+bX(P)+c_0)}=0,
\]
so \(X(P)^2=c_0\). Scaling by the nonzero \(X(P)\) normalizes the cubic to \(W^2=M(M^2+\alpha M+1)\), with \(M(P)=1\). A subgroup-preserving isomorphism fixes infinity, the selected two-torsion point at zero, and the common coordinate one of the two generators. The induced affine change of the cubic coordinate is therefore the identity, so it fixes \(\alpha\). Conversely \(\alpha\ne\pm2\) defines such a pair. Thus \(\nu=(\alpha+2)/4\) is a complete coordinate on the noncuspidal moduli problem, which is precisely \(Y_0(4)\)~\cite[\S19.3]{Sutherland}.

A generic elliptic curve has twelve exact-order-four points, forming six cyclic subgroups. Its only origin-preserving automorphisms are \(\pm1\), which preserve each subgroup. Hence forgetting the subgroup is generically six-to-one, in agreement with the degree of \eqref{eq:jnu}. Adding the three missing parameter values compactifies to a projective line, establishing the Hauptmodul statement.
\end{proof}

For example \(j(1-\nu)=j(\nu)\), corresponding to \(t\mapsto1/t\), but that identity alone does not describe the entire generic sixfold fiber. These are equivalences of bare complex curves, not equivalences of the fully labeled Three-Way cover. Neither do they identify real twists or fixed real plots.

The natural quotient by \(2P=(0,0)\) is explicit. Put
\[
M'=M+\alpha+M^{-1},\qquad W'=W(1-M^{-2}).
\]
Squaring and using \eqref{eq:Weierstrass} gives
\begin{equation}\label{eq:isogenous}
E':\qquad W'^2=M'(M'-(\alpha+2))(M'-(\alpha-2)).
\end{equation}
The map has degree two and kernel \(\{O,2P\}\). The nonzero roots are \(4\nu,4(\nu-1)\). Under \(s'=\nu-M'/4\), the three roots become \(\nu,0,1\); a nonzero constant rescaling of the ordinate therefore gives a Legendre equation with parameter \(\nu\). Thus
\begin{equation}\label{eq:jprime}
j(E')=256\frac{(1-\nu+\nu^2)^3}{\nu^2(1-\nu)^2}.
\end{equation}
This generally differs from \eqref{eq:jnu}. The parameter \(\nu\) is Legendre for the natural two-isogenous curve, whereas the displayed Legendre parameter for \(E_t\) itself is \(m\).

Let \(\overline E=E_t/H_{\rm tr}\). Since
\(H_{\rm tr}=[2]^{-1}(\langle2P\rangle)\), this quotient is also isomorphic to \(E'\). The map \(E_t\to\overline E\) is a degree-eight isogeny, realized by multiplication by two followed by the two-isogeny. The remaining quotient by inversion is a degree-two map \(\overline E\to\PP^1_\beta\) branched at \(0,1,\infty,\nu\). This explains geometrically why those four target points encode a Legendre curve different from the original closure.

\section{The marking lambda upstairs}\label{sec:divisor}

Let \(\pi:E_t\to\PP^1_\beta\) denote the degree-sixteen Galois map. For a regular marked value,
\begin{equation}\label{eq:Dlambda}
D_\lambda=\pi^*[\lambda],\qquad \lambda\notin\{0,1,\infty,\nu\},
\end{equation}
is a reduced effective divisor of degree sixteen. A Galois group acts transitively on every unramified fiber and the stabilizer is trivial, so this is one full free \(G\)-orbit. Choosing a lift \(R\) describes its support as
\begin{equation}\label{eq:orbitdivisor}
(R+H_{\rm tr})\ \cup\ (-R+H_{\rm tr}).
\end{equation}
On \(\overline E\), this is simply an unordered pair \(\{q,-q\}\). Neither one of the two points nor a particular lift is selected. Generically the points are not torsion.

The divisor class does not vary:
\[
D_\lambda\sim D_\infty,
\]
because their difference is the principal divisor of \(\beta-\lambda\). Hence \(\lambda\) specifies a particular effective divisor in this fiber family, not a new class in \(\operatorname{Pic}^{16}(E_t)\), and not another elliptic modulus. Under the \(F\)-decoration the four zeros and four poles on the \(x\)-sphere each lift to two points, splitting \(D_\lambda\) into eight zero points and eight pole points. That partition is extra labeling of the orbit, already accounted for in the automorphism discussion.

The two-modulus object can therefore be described either as the selected decorated branched cover over five labeled target points, or as the elliptic closure with its covering structure and a marked quotient orbit. It should not be replaced by an arbitrary elliptic curve with one marked point: that would forget the cyclic subgroup, the cover labels, and the unordered-pair structure.

\section{Degeneration, enhanced automorphisms, and topology}\label{sec:topology}

\subsection{The explicit limiting elliptic fibers}
The biquadratic equation gives concrete limits, including the components at infinity. At \(t=0\), it is \(x^2+z^2=0\), two rational \((1,1)\) components \(z=\pm ix\), meeting at \((0,0)\) and \((\infty,\infty)\). At \(t=\infty\), first divide the equation by \(t\); the limit is \(x^2z^2+1=0\), two components \(xz=\pm i\), meeting at \((0,\infty)\) and \((\infty,0)\). At \(t=1\), the factorization is \((x^2-1)(z^2-1)=0\); at \(t=-1\), it is \((x^2+1)(z^2+1)=0\), up to a nonzero sign. Each consists of two vertical and two horizontal rational components in a four-cycle.

The intersections are ordinary nodes. Differentiation with respect to \(t\), or \(1/t\) at infinity, shows that the nearby total family is smooth at these nodes. Thus this explicit semistable family over the \(t\)-line has polygon fibers \(I_2,I_2,I_4,I_4\). With only an elliptic origin retained, stabilization contracts unmarked rational components to a nodal rational stable genus-one curve. None is a smooth torus at the boundary.

The invariant has asymptotics
\[
j(\nu)\sim-16/\nu\quad(\nu\to0),\qquad
j(\nu)\sim16/(\nu-1)\quad(\nu\to1),\qquad
j(\nu)\sim65536\nu^4\quad(\nu\to\infty).
\]
The pole orders \(1,1,4\) are the cusp widths for this \(X_0(4)\) coordinate. The double covering by \(t\) ramifies at zero and infinity, explaining the two order-two poles there. A \(j\)-pole alone does not determine the fiber type of every twisted Weierstrass model; the component assertions above concern the explicit biquadratic family.

\subsection{Smooth extra automorphisms are different from boundaries}
For an elliptic curve with origin, the short equation \eqref{eq:short} shows that an automorphism has scaling form \((\xi,W)\mapsto(\omega^2\xi,\omega^3W)\), with \(\omega^4p=p\), \(\omega^6q=q\). Therefore \(\Aut(E,O)\) is \(C_2\) generically, \(C_4\) at \(j=1728\), and \(C_6\) at \(j=0\). Translations are additional automorphisms if no origin is fixed.

The exact smooth loci in the parameter line are
\begin{align*}
j=0:&\quad t^2=-7\pm4\sqrt3,
\qquad\nu=(2\pm\sqrt3)/4,\\
j=1728:&\quad t=\pm i\ \text{or}\ t^2\pm6t+1=0,\\
&\quad\nu=1/2\ \text{or}\ \nu=(4\pm3\sqrt2)/8.
\end{align*}
For the second statement one can check directly
\[
j(t)-1728=
\frac{16(t^2+1)^2(t^2-6t+1)^2(t^2+6t+1)^2}{t^2(t^2-1)^4}.
\]
In coefficient coordinates the corresponding equations are
\[
j=0:\ B^4+14B^2c^2+c^4=0,
\]
\[
j=1728:\ (B^2+c^2)(B^2-6Bc+c^2)(B^2+6Bc+c^2)=0,
\]
away from the excluded singular loci. For real \(t\), \(j=0\) never occurs; the real \(j=1728\) parameters are \(\pm(3\pm2\sqrt2)\).

These are smooth complex-multiplication points, not branch collisions. Extra units do not preserve the selected cyclic order-four subgroup: if such a unit acted as \(\pm1\) on a generator, an endomorphism with norm at most three would annihilate an order-four point, impossible since its kernel would have to contain four points. The deck group of the labeled degree-sixteen cover stays of order sixteen. Most CM curves have only the generic two origin-preserving automorphisms, so enhanced finite symmetry and CM are also different questions.

\subsection{The original source is a sphere; the closure is a torus}
The distinction is exact:
\begin{equation}\label{eq:spheretorus}
\boxed{\text{Original complex source }=\PP^1_x,\quad g=0,}
\qquad
\boxed{\text{Generic normal closure }=E_t,\quad g=1.}
\end{equation}
A smooth compact complex genus-one curve is analytically \(\CC/\Lambda\), hence topologically a torus. Its real locus, when defined over \(\RR\), consists of circles or is empty depending on the real form; it is not itself a real two-dimensional torus. The graph \(\{(x,f(x)):x\in\RR, D(x^2)\ne0\}\) is a collection of one-dimensional real arcs. Neither that graph nor the original sphere becomes a torus merely because its normal closure has genus one.

\fig{10_sphere_torus_distinction}{Two different operations: restrict a rational sphere map to a real graph, or pass to its complex normal closure and uniformize. The right-hand torus belongs to the latter operation.}{fig:topology}

A surface of revolution requires a specified planar profile, an axis, and choices at endpoints and singularities. For instance, rotating the full graph about the vertical ordinate axis produces the parametrized set
\[
(r\cos\theta,F(r^2),r\sin\theta),\qquad r\ge0,
\]
on each pole-free radial interval. The two signs of the original horizontal coordinate then describe the same circle, so \(x=\pm1\) become one radius-one circle at height one, not one point. A connected such branch is a radial graph: it has no self-intersection and, if it includes the axis regularly, a disk-like cap; otherwise it is annular. Its noncompact ends or pole-separated components do not create a compact torus. Adding the reciprocal profile can create intersections along common circles; reflecting about the lock line adds \(2-F(r^2)\), another prescribed profile. Additional cutting and gluing could define other objects, but the rational-map data alone do not select those operations.

The earlier requests to identify \(-1\) and \(1\) about a center, and to mirror across height one, are therefore compatible with discussing a separately prescribed construction, but are not proofs of its topology. The phase lift in H4 is likewise an added map with its own definitions. The appearance of an associated elliptic torus makes the historical intuition interesting in retrospect; it supplies no evidence for a toroidal universe, brane, or spacetime.

\section{Four historical coefficient pairs}\label{sec:examples}

All four documented examples have \(A>1\), \(c>0\), and \(\Delta>0\). We choose \(d=\sqrt\Delta>0\), so \(\rho=(c+B)/d>0\) and \(k>0\). Let \(\phi=(1+\sqrt5)/2\), and abbreviate \(J(t)\) by the rational expression in Eq.~\eqref{eq:jt}.

\begin{table}[htbp]\centering\small
\begin{tabular}{lllll}\toprule
Pair & \(\Delta\) & \(\rho\) & \(k\) & \(\lambda\)\\\midrule
\((9,9)\) & \(19\) & \(\sqrt{19}\) & \(16/\sqrt{19}\) & \(64/19\)\\
\((9,1)\) & \(99\) & \(\sqrt{11}/3\) & \(16/(3\sqrt{11})\) & \(64/99\)\\
\((\pi,e)\) & \((\pi+1)^2-e^2\) & \(\sqrt{\frac{\pi+1+e}{\pi+1-e}}\) & \(\frac{2(\pi-1)}{\sqrt\Delta}\) & \(\frac{(\pi-1)^2}{\Delta}\)\\
\((\phi,1-\phi)\) & \(4\phi\) & \(\phi^{-1/2}\) & \(\phi^{-3/2}\) & \((\sqrt5-2)/4\)\\\bottomrule
\end{tabular}\caption{Exact reduced-quotient data. In the \((\pi,e)\) row, \(\Delta\) means that row's exact value.}\label{tab:examples1}
\end{table}

\begin{table}[htbp]\centering\small
\begin{tabular}{llll}\toprule
Pair & \(t\) & \(\nu\) & \(m\)\\\midrule
\((9,9)\) & \(9/10\) & \(100/19\) & \(1/361\)\\
\((9,1)\) & \(1/10\) & \(100/99\) & \(81/121\)\\
\((\pi,e)\) & \(e/(\pi+1)\) & \(\frac{(\pi+1)^2}{(\pi+1)^2-e^2}\) & \(\frac{(\pi+1-e)^2}{(\pi+1+e)^2}\)\\
\((\phi,1-\phi)\) & \(2-\sqrt5\) & \((2+\sqrt5)/4\) & \(\phi^2\)\\\bottomrule
\end{tabular}\caption{Exact full-cover and elliptic data. The parameter \(m\) refers to the displayed Legendre model of the closure, not its two-isogenous quotient.}\label{tab:examples2}
\end{table}

\subsection{The pair (9,9)}
Here \(\chi=45\), \(\lambda=64/19>1\), and the marked fiber lies inside the real quotient interval. Its roots are
\[
y_z=(3\pm\sqrt5)/6,\qquad y_p=(9\pm3\sqrt5)/2.
\]
They are positive and distinct, producing four real zeros and four real poles in \(x\). This accounts for the archived asymptote list without an appeal to numerical symbolism. The critical \(y\)-pair is \((10\pm\sqrt{19})/9\), also positive. Both \(-1\) level points are nonreal because \(\Delta>0\).

The closure has
\[
j=\frac{8780093172522724}{263900025}
\simeq33270528.0059095258.
\]
This rational number is not an integer. Every CM \(j\)-invariant is an algebraic integer~\cite[Theorem 2.1]{DLR}, and a rational algebraic integer is an integer. Therefore this elliptic curve is non-CM. A large numerical value of \(j\) is not the proof; its exact arithmetic is.

\subsection{The pair (9,1)}
Here \(\chi=-35\) and \(0<\lambda=64/99<1\). The marked level is outside the real quotient image, so neither \(f\) nor its reciprocal has a real zero or pole. Nevertheless the exact real locks persist, and the positive critical pair \(10\pm\sqrt{99}\) produces a different smooth graph shape. Its \(-1\) fiber is again nonreal. The full cover has \(\nu=100/99\), close to the boundary value one but not on it.

The exact elliptic invariant is
\[
j=\frac{5927735656804}{2401490025}
\simeq2468.35739274161674,
\]
another rational noninteger, proving non-CM by the same integrality criterion. Thus the historical ambiguity between the two integer pairs changes both real geometry and elliptic covering data, despite their shared value of \(A\).

\subsection{The pair (pi,e)}
Elementary bounds imply \(e^2<4\pi\), hence \(\chi<0\) and \(0<\lambda<1\). This pair occupies the same zero/pole regime as \((9,1)\), with different values of both moduli. Its source scale exceeds one, so the critical points are on the positive \(y\)-half-line. The \(-1\) fiber is nonreal. At 100-digit working precision the principal invariants, rounded for display, are
\begin{align*}
\rho&\simeq2.195372442666,& k&\simeq1.370752047260,\\
\lambda&\simeq0.469740293767,&\nu&\simeq1.75678591761901002,\\
t&\simeq0.656337321369094558,&m&\simeq0.0430494062828726886,\\
j&\simeq132958.437104208270.&&
\end{align*}
Its exact \(j\) is \(J(e/(\pi+1))\). The CM status remains unresolved in this analysis. Transcendental input constants do not by themselves prove that a rational expression in them is transcendental. No further arithmetic investigation of this case is undertaken in this draft, and no proximity to any special numerical value is used as evidence.

\subsection{The pair (phi,1-phi)}
The golden-ratio relation \(\phi^2=\phi+1\) simplifies the exact invariants to those in Tables~\ref{tab:examples1}--\ref{tab:examples2}. Here \(B<0\), \(\rho<1\), and \(0<\lambda=(\sqrt5-2)/4<1\). There are no real zeros or poles and no nonzero real \(x\)-critical points: the real \(y\)-critical pair is negative. On \(y\ge0\), the function increases from \(1/\phi\) to \(\phi\), crossing height one exactly at \(y=1\). The \(-1\) level is nonreal. This additional embedding difference is not captured merely by saying that the marked fiber lies below one.

For \(m=\phi^2\), one has \(m^2-3m+1=0\), so \((m-1)^2=m\) and the Legendre formula reduces exactly to \(j=2048\). The complete list of thirteen rational CM \(j\)-values does not contain 2048~\cite[Introduction]{DLR}; consequently this curve is non-CM. Its algebraic simplification is real, but does not make it an enhanced-automorphism or CM point.

All four curves are smooth and none has \(j=0\) or \(1728\). For the transcendental pair, the elementary bound \(0.65<t<0.66\), compared with the explicit exceptional \(t\)-values in Section~\ref{sec:topology}, suffices to establish that statement without deciding CM.

\section{A bounded nearest-neighbor spectral comparison}\label{sec:physics}

This section begins with an independently specified standard chain and derives its objects before making a comparison. We do not change its Hamiltonian, fit its parameters to a graph, or postulate a new observable.

Let the uniform nearest-neighbor Hamiltonian on \(\ell^2(\ZZ)\) be
\begin{equation}\label{eq:chain}
(H\psi)_n=E_0\psi_n-J(\psi_{n+1}+\psi_{n-1}),\qquad J>0.
\end{equation}
The eigenvalue equation, with \(\epsilon=(E_0-E)/J\), becomes
\(\psi_{n+1}+\psi_{n-1}=\epsilon\psi_n\). Its Bloch ansatz \(\psi_n=u^n\) therefore yields
\begin{equation}\label{eq:dispersion}
\epsilon=u+u^{-1},\qquad E=E_0-J(u+u^{-1}).
\end{equation}
This is the familiar tight-binding dispersion derived, for example, in Tong~\cite[\S2.1.1, Eqs.~(2.1)--(2.5)]{Tong}. Quotienting the two multipliers by exchange \(u\leftrightarrow u^{-1}\), and then identifying energies relative to the band center by \(\epsilon\leftrightarrow-\epsilon\), gives
\begin{equation}\label{eq:spectralquotient}
\beta=\epsilon^2/4.
\end{equation}
This is exactly the reduced quotient already derived in Section~\ref{sec:belyi}, for every generic coefficient pair. The chain has not acquired new couplings depending on \(A,B\); those coefficients select a source coordinate and marked spectral value \(\lambda=k^2/4\).

\subsection{Domains cannot be identified by the formula alone}
Propagating real-energy states have \(u=e^{iq}\), real \(q\), so \(|u|=1\), \(\epsilon=2\cos q\in[-2,2]\), and \(\beta\in[0,1]\). For real \(\epsilon\) outside the band, the two multipliers are real and reciprocal; one grows and one decays. A decaying multiplier can enter a resolvent or a boundary problem, but it is not a propagating Bloch wave. At the band edges the multipliers coalesce.

For the original real plotting contour with \(\Delta>0\), \(u\) is real, not generally on the unit circle; it samples this evanescent algebraic continuation. Its marked \(k\) is real, with \(0<\lambda<1\), \(\lambda=1\), and \(\lambda>1\) corresponding to a level inside the band, at an edge, or outside it. A marked level can lie inside the band even when the original real plotting contour does not reach it. For \(\Delta<0\), both \(u\) and \(k\) are imaginary on the chosen real source contour, and \(\epsilon\) is generally imaginary. That is a complex spectral continuation, not a real-energy evanescent state. These restrictions survive the full two-parameter normalization.

\subsection{Standard chain objects and the requested ratio}
The transfer matrix derived from the recurrence is
\[
\begin{pmatrix}\psi_{n+1}\\\psi_n\end{pmatrix}
=\begin{pmatrix}\epsilon&-1\\1&0\end{pmatrix}
\begin{pmatrix}\psi_n\\\psi_{n-1}\end{pmatrix}.
\]
Its eigenvalues are \(u,u^{-1}\); the adjacent-amplitude ratio of a Bloch solution is \(u\). These quantities contain the characteristic square root \(\sqrt{\epsilon^2-4}\). No independent level \(k\) appears.

To fix resolvent signs, let \(\mathcal A=S+S^*\) be the dimensionless adjacency operator, and define Green functions with \((z-\mathcal A)^{-1}\). On the semi-infinite chain, a Schur complement at the endpoint gives
\[
g(z)=\frac1{z-g(z)},\qquad
g(z)=\frac{z-\sqrt{z^2-4}}2,
\]
where the branch is chosen so that \(g(z)\sim z^{-1}\) at infinity. The alternative Weyl convention \(\langle\delta_0,(\mathcal A-z)^{-1}\delta_0\rangle\) is \(-g\). On the full line the diagonal Green function is \(1/\sqrt{z^2-4}\), as follows either by Fourier integration or by joining two half-line self-energies:
\(1/(z-2g(z))=1/\sqrt{z^2-4}\). A hopping self-energy is the appropriately dimensioned multiple of the half-line function. These independently defined functions are not the desired rational ratio.

For a uniform full chain there is no scattering defect: reflection is zero and transmission is a reference-dependent propagation phase. For a hard endpoint with \(\psi_0=0\), writing \(\psi_n=e^{-iqn}+r e^{iqn}\) gives \(r=-1\). Moving the reference plane only multiplies by a known phase. A more general endpoint impedance could introduce a boundary parameter, but selecting one to force a desired formula would change the stated problem. These standard cases therefore do not independently yield \((\epsilon+k)/(\epsilon-k)\).

There is a precise but limited Green-function identity. Fourier diagonalization gives the momentum-resolved resolvent
\[
G(z;q)=\frac1{z-2\cos q}.
\]
Choose a momentum \(q_0\) and write \(k=2\cos q_0\). The opposite band-center level occurs at \(\pi-q_0\), so
\begin{equation}\label{eq:Greenratio}
\frac{G(\epsilon;q_0)}{G(\epsilon;\pi-q_0)}
=\frac{\epsilon+k}{\epsilon-k}=F.
\end{equation}
The Green functions are independently defined standard objects; the choice of two momenta and the decision to take this ratio are not selected by the uniform chain itself. Thus \eqref{eq:Greenratio} is an exact constructed comparison, not an independently established observable dictionary. For \(|k|>2\), \(q_0\) is complex and even the selected momentum is outside the real Bloch band; imaginary \(k\) requires complex continuation as well.

The limited negative result is therefore that no independently selected chain observable equal to the full nonconstant \(F\) has been established among these objects. It is not a theorem that no conceivable reformulation could ever contain this ratio. That stronger claim would require a defined class of all observables and boundary problems, which is not part of this task.

\subsection{Does the additional square root have a physical meaning?}
The chain naturally has a two-sheeted multiplier relation: solving \(u+u^{-1}=\epsilon\) interchanges the growing and decaying, or left- and right-moving, branches. That is already the Joukowski part of the quotient. Spatial reflection exchanges \(u\) and \(u^{-1}\); staggering \(\psi_n\mapsto(-1)^n\psi_n\) exchanges \(u\) and \(-u\) and reverses energy about the band center. These explain the reduced algebraic symmetries.

They do not independently produce
\[
y=x^2,\qquad u=\rho\frac{y-1}{y+1}.
\]
The additional branching at \(u=\pm\rho\), with image \(\nu\), is selected by the rational construction and its coordinate embedding, not by a standard chain parity reduction or an amplitude-squaring rule. Calling \(x\) a square root of a multiplier would ignore the intervening Möbius transformation; calling it an amplitude would define a meaning rather than derive one. No independent physical meaning for the extra \(x\)-cover was found. The successful dictionary stops at the reduced complex spectral quotient \eqref{eq:spectralquotient}.

\section{Historical interpretation and what the reconstruction establishes}\label{sec:lessons}

The historical papers used toroidal, brane, inter-universal, and gravitational language to motivate or interpret exploratory geometry~\cite{H1,H2,H4}. Those statements remain part of the historical record. They are not assumptions needed for the mathematical proofs, and the present reconstruction supplies neither dynamics nor a metric, stress tensor, matter model, empirical coupling, or predictive identification that would establish them as physics. Retaining that distinction respects the historical exploration without preserving unsupported conclusions.

There are three kinds of results here. Exact identities and classifications follow from the displayed algebra and covering arguments. Numerical checks provide bounded independent consistency tests of those identities and limits. Graphs and topological schematics make their geometric interpretation visible, but do not replace the proofs. The later reconstruction is more precise than the early imagery because it specifies which object carries each claimed property.

Many individual structures are explicitly identified with established mathematics. Basic reciprocity, parity, and intersections were also already present in historical or local analytical material. The manuscript's reconstruction includes the coherent organization of those facts, the additional involution, the marked-fiber classification, the multiscale chart relations, and the exact normal-closure identification. Whether that particular assembly or any specific formulation has appeared elsewhere is deliberately left to a dedicated prior-art audit. Neither the use of special constants nor an unexpectedly simple invariant is evidence of novelty.

The remaining questions have narrower scopes than the early physical claims. A future prior-art audit can compare the precise tower and parameter maps with existing cover and elliptic constructions. A separate geometrical work order could specify a complete rotation and gluing rule for the visual object. A physical investigation would first need an independently defined observable and a justified meaning for the extra source cover. The CM status of the \((\pi,e)\) curve remains recorded as unresolved, without further investigation here.

\begin{longtable}{L{51mm}L{101mm}}
\caption{Historical questions and modern answers.}\label{tab:lessons}\\
\toprule Historical question & Modern answer\\\midrule\endfirsthead
\toprule Historical question & Modern answer\\\midrule\endhead
Is reciprocity real? & Yes. Coefficient reversal gives an exact rational identity, with poles, zeros, and cancellation handled projectively.\\
Why are \(x=\pm1\) special? & They give \(y=1\), where numerator and denominator agree. The height is exactly one unless the expression cancels or is undefined; cancellation must be analyzed separately.\\
Do the historical real plots also reach \(-1\)? & No for all four examples studied here: \(\Delta>0\) makes that fiber nonreal. The complex \(-1\) fiber still exists.\\
What controls zeros and poles? & \(\chi=B^2-4A\) controls repeated roots and real-root transitions; \(\Delta=0\) and \(A=1\) control cancellation. The positive \(y\)-restriction decides which roots are visible in real \(x\).\\
Was the number \(9\) the mechanism? & No. The algebraic symmetries hold throughout the generic family. \((9,9)\) and \((9,1)\) occupy different strata and cannot be merged.\\
Why do large coefficients create nested scales? & Distinct balances resolve an inner zero, an outer pole, a thin lock, and intermediate stationary shoulders. Exact overlap identities prove the matching.\\
Where does \(s\mapsto2/s\) come from? & It is the scaled limit of the equal-value involution and exchanges the two matched contributions. Reciprocal inversion instead tends to \(s\mapsto-s\).\\
What are \(\lambda\) and \(\nu\)? & \(\lambda\) marks the zero/pole fiber; \(\nu\) is the additional branch value from squaring and the modular coordinate for the selected elliptic order-four structure.\\
Is the full object a five-marked sphere? & Its selected generic decorated coarse space is \(M_{0,5}\). The Hurwitz object also retains a covering type and nontrivial automorphisms.\\
Is the construction a sphere or a torus? & The original complex source is a sphere; its generic normal closure is genus one and analytically a torus. A revolved real profile is a separate construction.\\
Is there established physics? & There is an exact reduced nearest-neighbor spectral-quotient match. No full physical theory follows from it.\\
What remains physically unexplained? & No independently selected observable for \(F\), and no independent physical role for the additional \(x\)-cover, has been established.\\
\bottomrule
\end{longtable}

\appendix
\section{Evidence architecture and reproducibility}\label{app:evidence}

This edition preserves five external analytical reports dated 5 October 2026: the large-\(L\) report R1~\cite{R1}, known-structure review R2~\cite{R2}, full two-parameter normal form R3~\cite{R3}, full-cover moduli report R4~\cite{R4}, and elliptic closure report R5~\cite{R5}. Their original files remain unchanged. The snapshots retain their original wording, including statements subsequently refined by another report. They are evidence of the reconstruction sequence, not substitutes for the final arguments here. In particular the earlier physical-analogue review's open follow-up is resolved only to the bounded negative extent stated in Section~\ref{sec:physics}.

The paper is built directly from the editable LaTeX source, with vector mathematical figures generated by a deterministic script. The supplementary directory contains the bibliography in BibTeX and rendered-reference forms, a historical source index, the complete machine-readable claim ledger, exact source hashes, verification code and results, and a build record. The repository baseline is recorded before authoring. Files in the scientific kernel, UI, and historical corpus are hashed for before/after comparison. No commit or publication operation is part of this edition.

\subsection{Independent computational checks}
The verification script evaluates 77 exact symbolic identities with SymPy. They include reciprocity, the historical \(C\)-extension, level-set differences, resultants, derivatives, parameter transformations, the coherent normal form, the marked-fiber identity, intermediate-chart actions, scaling and quotient-chart limits, elliptic transformations, discriminants, the two-isogeny, enhanced-locus factorization, and exact historical \(j\)-values. Rational identities are reduced algebraically; numerical tolerance is not used to declare them true.

The branch-cycle implementation explicitly generates the permutation group from Eq.~\eqref{eq:cycles}, confirming transitivity, order sixteen, and identity product. This checks the finite presentation used in the proof, rather than searching for a group from approximate data.

At 100-decimal-digit working precision, twelve fixed parameter pairs and four fixed nonreal source points per pair give 48 normal-form comparisons. They include positive and negative \(\Delta\), negative \(A\), \(A=0\), \(A=-1\), a constant-map case, both signs of \(B\), and a case close to cancellation. The maximum normalized discrepancy is \(2.315\times10^{-100}\). These points are deliberately off the denominator zeros and do not assert numerical stability at a pole. The four historical examples are evaluated at the same precision; their three formulas for \(j\) agree to working precision. The displayed decimals in the body are rounded, whereas the supplementary JSON retains substantially more digits.

For the intermediate limit we use the six fixed points \(s=\pm1/2,\pm\sqrt2,\pm3\). The following table reports absolute profile errors, together with relative errors for two matching paths \(\delta=L^{-1/4}\) (bulk side) and \(\delta=L^{-3/4}\) (lock side). The exact uniform estimate in Theorem~\ref{thm:matching}, not the table, proves matching.
\begin{table}[htbp]\centering\small
\begin{tabular}{rrrr}\toprule
\(L\) & \(\max|P_h-W|\) & Bulk relative error & Lock relative error\\\midrule
\(10^4\) & \(1.1324\times10^{-1}\) & \(2.3338\times10^{-2}\) & \(1.1803\times10^{-1}\)\\
\(10^8\) & \(1.1003\times10^{-3}\) & \(2.0303\times10^{-4}\) & \(1.0153\times10^{-2}\)\\
\(10^{12}\) & \(1.1000\times10^{-5}\) & \(2.0030\times10^{-6}\) & \(1.0015\times10^{-3}\)\\\bottomrule
\end{tabular}\caption{Bounded high-precision consistency checks of the proved scaling formulas. Relative errors compare \(f_L+1\) with \(-\delta\) and \(-2/(L\delta)\), respectively.}\label{tab:numerics}
\end{table}

\subsection{Claim ledger summary}
The full ledger assigns each claim a proof location, an analytical report, a historical witness where applicable, a literature anchor where used, and an evidence status. The following summary makes the distinction visible in the manuscript itself. A dash in the history column means the modern claim is not attributed to the historical document.

\begin{longtable}{L{13mm}L{48mm}L{39mm}L{45mm}}
\caption{Evidence and proof map. ``Exact'' refers to the proof in this manuscript; historical statements are documentary assertions about a source.}\label{tab:ledger}\\
\toprule ID & Claim & Manuscript proof & Evidence context\\\midrule\endfirsthead
\toprule ID & Claim & Manuscript proof & Evidence context\\\midrule\endhead
C01 & Minus-sign target; conflicting labels and later plus sign & \S\ref{sec:history} (documentary) & H1 pp.~13--14; H2 PDF, code and CSVs; H3 p.~3; H4 p.~1.\\
C02 & Reciprocity and lock fibers & Prop.~\ref{prop:levels} & H1, H3; R3; M1. Exact algebra.\\
C03 & Roots, cancellations, and real regimes & Prop.~\ref{prop:divisors}; \S\ref{sec:elementary} & R3; corrects H1 p.~14. Exact algebra.\\
C04 & Equal-value involution and coherent normalization & Thm.~\ref{thm:normal} & R1, R3. Exact algebra.\\
C05 & Universal labeled \(V_4\) quotient & Thm.~\ref{thm:belyi} & R1--R3; standard Belyi/orbifold species~\cite{SV,KO}.\\
C06 & Completeness of reduced \(\lambda\) & \S\ref{sec:belyi} & R3. Equivalence relation explicitly fixed.\\
C07 & Real contour differs from complex equivalence & \S\ref{sec:real} & R3, R4. Exact domain restrictions.\\
C08 & Large-\(L\) charts and matching & Thm.~\ref{thm:matching}; \S\ref{sec:scaling} & R1. Proof plus bounded numeric checks.\\
C09 & Distinct scaled origins of the involutions & Eqs.~\eqref{eq:Rh}--\eqref{eq:shoulders} & R1 follow-ups. Exact finite-scale identities.\\
C10 & Degree-eight branch data and lift exceptions & Thm.~\ref{thm:xcover}; \S\ref{sec:xcover} & R3, R4. Valuation and ramification proofs.\\
C11 & Selected coarse moduli \(M_{0,5}\), four presentations & Thm.~\ref{thm:moduli} & R4; marked-target convention~\cite{Mnev}.\\
C12 & Cover type and stack qualification & \S\ref{sec:moduli}; \S\ref{sec:galois} & R4. Pair types and intrinsic intermediate.\\
C13 & Normal closure, genus and monodromy & Thm.~\ref{thm:galois}; Eq.~\eqref{eq:cycles} & R4, R5. Exact function-field and group arguments.\\
C14 & Stable-target and admissible-cover boundaries & \S\ref{sec:boundary} & R4; standard framework~\cite{Braungardt,CMR}.\\
C15 & Edwards, Montgomery, Jacobi, Legendre models & \S\ref{sec:edwards} & R5; named models~\cite{Edwards,Morain,EFD}.\\
C16 & Elliptic group action and two-torsion quotients & \S\ref{sec:group} & R5; addition law~\cite{EFD}.\\
C17 & \(\nu\) as \(X_0(4)\) coordinate; two-isogeny & Thm.~\ref{thm:X04}; Eq.~\eqref{eq:isogenous} & R5; modular interpretation~\cite{Sutherland,Morain}.\\
C18 & Lifted \(\lambda\)-orbit and divisor class & \S\ref{sec:divisor} & R5. Exact quotient/divisor argument.\\
C19 & Sphere, elliptic torus, and revolution differ & \S\ref{sec:topology} & R4, R5; historical H1/H4 interpretation delimited.\\
C20 & Historical invariants and CM status & \S\ref{sec:examples} & R3, R5; exact CM criteria~\cite{DLR}. \((\pi,e)\) unresolved.\\
C21 & Reduced spectral comparison and bounded negative dictionary & \S\ref{sec:physics} & R2, R3; chain derivation~\cite{Tong}. No full observable established.\\
C22 & No priority claim for the assembled reconstruction & \S\ref{sec:intro}; \S\ref{sec:lessons} & Dedicated prior-art audit pending.\\
\bottomrule
\end{longtable}

\section{Limits of this edition and exact claims for a subsequent audit}\label{app:audit}

The documentary record leaves the intended meaning of some historical labels unresolved; both formulas are retained rather than harmonized. Printed dates have been checked against the archived PDFs, but live Zenodo deposit metadata was not verified in this edition. The historical stability and curvature experiments are indexed, not independently rerun. There is no broad novelty search, no arithmetic extension of the unresolved \((\pi,e)\) CM question, and no physical model added to the uniform chain.

The next prior-art audit should test the following precisely delimited statements, rather than asking whether reciprocity or elliptic curves in general are new:
\begin{enumerate}
\item The exact embedding of this reversed-quadratic family into \(F=(r+k)/(r-k)\), with \(u=\rho(y-1)/(y+1)\), its equal-value involution, and the interpretation of \(\lambda\) under the specified labeled equivalence.
\item The full squaring tower \(Q_t\), its \(R\)-paired Hurwitz type, the coarse \(M_{0,5}\) classification, and the four finite coefficient presentations.
\item The explicit genus-one normal closure \(tx^2z^2-x^2-z^2+t=0\), its degree-sixteen labeled quotient, and the interpretation of the two sphere-level \(V_4\) actions through two-isogenous Kummer quotients.
\item The particular parameter dictionary \(t=B/(A+1)\), \(\nu=(1-t^2)^{-1}\), the selected cyclic order-four subgroup, the two-isogenous Legendre parameter, and the lifted regular-fiber divisor \(D_\lambda\).
\item The diagonal multiscale reconstruction, especially the rigorously matched \(W(s)=-s-2/s\) profile and the exact finite-\(L\) origins of its limiting involutions.
\item The historical provenance reconciliation and correction of pole/lock claims, which is a documentary contribution distinct from mathematical priority.
\end{enumerate}
The audit must compare the same decorations and equivalences. A known elliptic equation can settle a model identification while leaving a particular coordinate tower to be traced. Conversely, a change of variables may show that the whole formulation is standard. The outcome of that audit remains open.

\clearpage
\phantomsection\addcontentsline{toc}{section}{References}
\begin{thebibliography}{99}
\bibitem{H1} H. F. Halldórsson and ChatGPT o4-mini-high (printed attribution), \emph{The Three-Way Quantum Gravitational Wave Model: A Refined Geometric and Quantum Field Theoretic Framework}, archived final draft, printed 1 May 2025, 14 pp. Historical source H1, especially Appendix C, pp.~13--14. Exact PDF and hash in the source manifest; corpus-associated Zenodo record 17240087.
\bibitem{H2} H. F. Halldórsson and ChatGPT-5 (printed attribution), \emph{Bounded Infinity and Resonant Curvature: A Geometric Analysis of Reciprocal Rational Maps in the Three-Way Resonance Model (RTM)}, printed 16 October 2025, 35 pp. Historical source H2, with archived \texttt{rrm\_tool.py} and CSVs; corpus-associated Zenodo record 17365334. Exact files and hashes in the source manifest.
\bibitem{H3} H. F. Halldórsson and ChatGPT 5.2 (printed attribution), \emph{Experimental Geometry of a Three-Way Reciprocal Construct: Numerical Evidence for Structural Stability}, printed 19 January 2026, 15 pp. Historical source H3, \texttt{Three\_way\_test\_model.pdf}; Eq.~(1), p.~3. No DOI assigned in this bibliography.
\bibitem{H4} H. F. Halldórsson and ChatGPT 5.4 (printed attribution), \emph{Reciprocal Rational Maps, Self-Dual Hinges, and Toroidal Phase Lifts}, printed 18 March 2026, 11 pp. Historical source H4; exact PDF and hash in the source manifest.
\bibitem{M1} \emph{Three-Way Item 2A: Constant Geometry; Item 2B: Constant Survey; Item 2C: M2 Fibonacci Structure}, local analytical notes, archived as M1 in this edition. These use the plus-sign convention and are not historical formula authority. Original file paths and hashes are retained in the source manifest.
\bibitem{R1} \emph{Three-Way Large-L Scaling Analysis v0.1}, external reconstruction report and bounded follow-ups, 5 October 2026. Evidence R1; includes shoulder, intermediate-involution, and quotient verification files. Verbatim snapshot supplied with this paper.
\bibitem{R2} \emph{Three-Way Known Structure and Physical Analogue Review v0.1}, external reconstruction report, 5 October 2026. Evidence R2; verbatim snapshot supplied with this paper.
\bibitem{R3} \emph{Three-Way Two-Parameter Normal Form v0.1}, external reconstruction report, 5 October 2026. Evidence R3; verbatim snapshot supplied with this paper.
\bibitem{R4} \emph{Three-Way X-Cover Moduli v0.1}, external reconstruction report, 5 October 2026. Evidence R4; verbatim snapshot supplied with this paper.
\bibitem{R5} \emph{Three-Way Genus-One Galois Closure v0.1}, external reconstruction report, 5 October 2026. Evidence R5; verbatim report, verification script, and results supplied with this paper.
\bibitem{SV} J. Sijsling and J. Voight, On computing Belyi maps, \emph{Publications mathématiques de Besançon. Algèbre et théorie des nombres} (2014), 73--131. \href{https://doi.org/10.5802/pmb.5}{doi:10.5802/pmb.5}. Accessible preprint: \url{https://arxiv.org/pdf/1311.2529}.
\bibitem{KO} J. Kalliongis and R. Ohashi, Finite actions on the Klein four-orbifold and prism manifolds, \emph{Commentationes Mathematicae Universitatis Carolinae} \textbf{58}(1) (2017), 49--68. \href{https://doi.org/10.14712/1213-7243.2015.193}{doi:10.14712/1213-7243.2015.193}. \url{https://cmuc.karlin.mff.cuni.cz/pdf/cmuc1701/kallio.pdf}.
\bibitem{Mnev} P. Mnev, \emph{Conformal Field Theory}, lecture notes, 2022, pp.~57--60, especially Eqs.~(209)--(210), for the marked-sphere configuration spaces. \url{https://academicweb.nd.edu/~pmnev/f22/CFT_2022_notes.pdf}.
\bibitem{Braungardt} V. Braungardt, Covers of moduli surfaces, \emph{Compositio Mathematica} \textbf{140}(4) (2004), 1033--1036, \S2. \href{https://doi.org/10.1112/S0010437X03000642}{doi:10.1112/S0010437X03000642}.
\bibitem{CMR} R. Cavalieri, H. Markwig, and D. Ranganathan, \emph{Tropicalizing the space of admissible covers}, arXiv:1401.4626, consulted preprint, Definition 8 and \S5.1. \url{https://arxiv.org/pdf/1401.4626}.
\bibitem{Edwards} H. M. Edwards, A normal form for elliptic curves, \emph{Bulletin of the American Mathematical Society} \textbf{44}(3) (2007), 393--422. \href{https://doi.org/10.1090/S0273-0979-07-01153-6}{doi:10.1090/S0273-0979-07-01153-6}. Publication information and accessible copy indexed by Bernstein and Lange: \url{https://www.hyperelliptic.org/tanja/newelliptic/newelliptic.html}.
\bibitem{Morain} F. Morain, \emph{Edwards curves and CM curves}, 2009, arXiv:0904.2243, especially \S\S2.4--2.5. \url{https://cacr.uwaterloo.ca/techreports/2009/cacr2009-17.pdf}.
\bibitem{EFD} D. J. Bernstein and T. Lange, \emph{Explicit-Formulas Database}, genus-one models, accessed 5 October 2026. Twisted Edwards: \url{https://www.hyperelliptic.org/EFD/g1p/auto-twisted.html}; Jacobi quartic: \url{https://www.hyperelliptic.org/EFD/g1p/auto-jquartic.html}; Montgomery: \url{https://www.hyperelliptic.org/EFD/g1p/auto-montgom.html}.
\bibitem{Sutherland} A. V. Sutherland, \emph{18.786 Lecture 19}, MIT, 2024, \S19.3 (p.~4), modular curves as moduli of elliptic curves with level structure. \url{https://math.mit.edu/classes/18.786/2024/LectureNotes19.pdf}.
\bibitem{DLR} H. B. Daniels and Á. Lozano-Robledo, \emph{On the number of isomorphism classes of CM elliptic curves defined over a number field}, author-hosted manuscript, Introduction and Theorem 2.1. \url{https://alozano.clas.uconn.edu/wp-content/uploads/sites/490/2014/01/Daniels_and_Lozano-Robledo_2.pdf}.
\bibitem{Tong} D. Tong, \emph{Lectures on Solid State Physics}, 2017, Chapter 2, \S2.1.1, pp.~27--29, Eqs.~(2.1)--(2.5). \url{https://davidtong.org/pdfs/teaching/solid-state-physics/solidstate2.pdf}.
\end{thebibliography}

\end{document}
